\documentclass[12pt, a4paper]{article}

% --- PAQUETES ---
\usepackage[utf8]{inputenc}
\usepackage[spanish]{babel}
\usepackage{amsmath}
\usepackage{amssymb}
\usepackage{amsfonts}
\usepackage{geometry}
\usepackage[most]{tcolorbox}
\usepackage{siunitx}
\usepackage{enumitem}
\usepackage{pgfplots}
\usepackage{tikz}
\usepackage{verbatim} 
\usepackage{xcolor}
\usepackage{listings} % Para mostrar código

\pgfplotsset{compat=1.18}
\usetikzlibrary{arrows.meta, calc, babel, angles, quotes}
\geometry{a4paper, margin=1in}

% --- CONFIGURACIÓN DE LISTINGS (PYTHON) ---
\definecolor{codegreen}{rgb}{0,0.6,0}
\definecolor{codegray}{rgb}{0.5,0.5,0.5}
\definecolor{codepurple}{rgb}{0.58,0,0.82}
\definecolor{backcolour}{rgb}{0.95,0.95,0.92}

\lstset{
    backgroundcolor=\color{backcolour},   
    commentstyle=\color{codegreen},
    keywordstyle=\color{magenta},
    numberstyle=\tiny\color{codegray},
    stringstyle=\color{codepurple},
    basicstyle=\ttfamily\footnotesize,
    breaklines=true,                 
    captionpos=b,                    
    keepspaces=true,                 
    numbers=left,                    
    numbersep=5pt,                  
    showspaces=false,                
    showstringspaces=false,
    showtabs=false,                  
    tabsize=2,
    language=Python
}

\title{Ejercicios: Movimiento Circular}
\author{Dataset Generado}
\date{\today}

\begin{document}

\maketitle
\subsection*{Enunciado}
Una rueda $A$ de radio $R = 0.4 \, \si{m}$ parte del reposo y aumenta su velocidad angular uniformemente a razón de $\alpha_A = 0.6\pi \, \si{rad/s^2}$. La rueda $A$ transmite su movimiento a una rueda $B$ mediante una correa. El radio de $B$ es $r = 0.15 \, \si{m}$.

¿Cuál es el tiempo necesario para que $B$ alcance una rapidez angular de $300 \, \text{rpm}$?

\begin{center}
\begin{tikzpicture}[scale=1.5]
    % Rueda A (Grande) - Motriz
    \draw[thick, fill=gray!10] (0,0) circle (0.8);
    \node at (0,0) {\textbf{A}};
    \draw[->] (0,0) -- (0,0.8) node[midway, left] {$R$};
    % Flecha de rotación A
    \draw[->, thick, blue] (0.3,0.3) arc (45:-45:0.4) node[midway, right] {$\alpha_A$};

    % Rueda B (Pequeña) - Conducida
    \draw[thick, fill=gray!10] (2.5,0) circle (0.3);
    \node at (2.5,0) {\textbf{B}};
    \draw[->] (2.5,0) -- (2.5,0.3) node[midway, left] {\tiny $r$};

    % Correa (Tangentes)
    \draw[thick] (0,0.8) -- (2.5,0.3);
    \draw[thick] (0,-0.8) -- (2.5,-0.3);

    % Etiquetas adicionales
    \node[below, gray] at (0,-0.9) {Motriz};
    \node[below, gray] at (2.5,-0.4) {Conducida};
\end{tikzpicture}
\end{center}

\subsection*{Solución}

\subsubsection*{1. Conversión de Unidades}
Primero, convertimos la velocidad angular deseada de la rueda B a unidades del Sistema Internacional (rad/s):
$$ \omega_{B_f} = 300 \, \text{rpm} \times \frac{2\pi \, \text{rad}}{60 \, \text{s}} = 10\pi \, \text{rad/s} \approx 31.42 \, \text{rad/s} $$

\subsubsection*{2. Principio de Transmisión}
Dado que las ruedas están unidas por una correa inextensible que no desliza, la \textbf{rapidez lineal (tangencial)} en el borde de ambas ruedas debe ser la misma ($v_A = v_B$).

Relacionamos la velocidad lineal con la angular ($v = \omega \cdot \text{radio}$):
$$ \omega_{A_f} R = \omega_{B_f} r $$

Despejamos la velocidad angular final necesaria en la rueda A:
$$ \omega_{A_f} = \omega_{B_f} \left( \frac{r}{R} \right) $$
$$ \omega_{A_f} = (10\pi) \left( \frac{0.15}{0.4} \right) = 3.75\pi \, \text{rad/s} \approx 11.78 \, \text{rad/s} $$

\subsubsection*{3. Cinemática Rotacional (Rueda A)}
Usamos la ecuación de velocidad para movimiento circular uniformemente variado (MCUV), partiendo del reposo ($\omega_{A_0} = 0$):
$$ \omega_{A_f} = \omega_{A_0} + \alpha_A t $$
$$ 3.75\pi = 0 + (0.6\pi) t $$

Despejamos el tiempo $t$. Note que $\pi$ se cancela en ambos lados:
$$ t = \frac{3.75\pi}{0.6\pi} $$
$$ t = \frac{3.75}{0.6} $$
$$ \boxed{t = 6.25 \, \text{s}} $$

\begin{tcolorbox}[colback=blue!5!white, colframe=blue!75!black, title=Respuesta Final]
\centering
El tiempo necesario es $t = 6.25 \, \text{s}$.
\end{tcolorbox}

\newpage

\subsection*{Enunciado}
Una partícula se mueve en una circunferencia de $0.7 \, \si{m}$ de radio, en el sentido contrario al avance de las manecillas del reloj con rapidez angular constante $\omega = 2 \, \si{rad/s}$.

A $t = 0$ parte del punto $A$ (ubicado a un ángulo de $\pi/6 \, \si{rad}$ respecto a la horizontal). Calcule el vector aceleración media para el intervalo de $t_1 = 1.7 \, \si{s}$ hasta $t_2 = 3 \, \si{s}$.

\begin{center}
\begin{tikzpicture}[scale=2, >=Latex]
    % Ejes
    \draw[->, gray] (-1.2,0) -- (1.2,0) node[right] {$x$};
    \draw[->, gray] (0,-1.2) -- (0,1.2) node[above] {$y$};
    
    % Círculo
    \draw[thick] (0,0) circle (1);
    
    % Radio y Punto A (t=0)
    \coordinate (O) at (0,0);
    \coordinate (A) at (30:1);
    
    \draw[dashed] (O) -- (A) node[midway, above left] {$R=0.7$};
    \fill[red] (A) circle (1.5pt) node[right] {$A (t=0)$};
    
    % Ángulo inicial
    \draw (0.3,0) arc (0:30:0.3);
    \node at (15:0.5) {\small $\pi/6$};
    
    % Flecha de rotación (omega)
    \draw[->, blue, thick] (60:1.2) arc (60:120:1.2) node[midway, above] {$\omega = 2$ rad/s};
\end{tikzpicture}
\end{center}

\subsection*{Solución}

\subsubsection*{1. Definición de Velocidad Lineal}
En un Movimiento Circular Uniforme (MCU), la magnitud de la velocidad es constante:
$$ v = \omega R = (2)(0.7) = 1.4 \, \si{m/s} $$

La aceleración media se define vectorialmante como:
$$ \vec{a}_m = \frac{\vec{v}_2 - \vec{v}_1}{t_2 - t_1} $$

Para hallar los vectores velocidad $\vec{v}_1$ y $\vec{v}_2$, primero necesitamos la posición angular de la partícula en cada instante. La ecuación de posición es $\theta(t) = \theta_0 + \omega t$, donde $\theta_0 = \pi/6$ rad.

\subsubsection*{2. Análisis para $t_1 = 1.7 \, \si{s}$}
Posición angular:
$$ \theta_1 = \frac{\pi}{6} + (2)(1.7) = 0.5236 + 3.4 = 3.9236 \, \si{rad} $$
Convirtiendo a grados: $3.9236 \times \frac{180}{\pi} \approx 224.8^\circ$ (III Cuadrante).

El vector velocidad es tangente a la trayectoria ($90^\circ$ adelante del radio). Usamos la definición vectorial:
$$ \vec{v}_1 = v (-\sin\theta_1 \hat{i} + \cos\theta_1 \hat{j}) $$
$$ \vec{v}_1 = 1.4 [-\sin(224.8^\circ)\hat{i} + \cos(224.8^\circ)\hat{j}] $$
$$ \vec{v}_1 = 1.4 [-(-0.704)\hat{i} + (-0.710)\hat{j}] $$
$$ \vec{v}_1 = 0.98\hat{i} - 0.99\hat{j} \, \si{m/s} $$

\subsubsection*{3. Análisis para $t_2 = 3 \, \si{s}$}
Posición angular:
$$ \theta_2 = \frac{\pi}{6} + (2)(3) = 0.5236 + 6 = 6.5236 \, \si{rad} $$
Convirtiendo a grados: $6.5236 \times \frac{180}{\pi} \approx 373.8^\circ$ (equivalente a $13.8^\circ$, I Cuadrante).

Vector velocidad:
$$ \vec{v}_2 = 1.4 [-\sin(373.8^\circ)\hat{i} + \cos(373.8^\circ)\hat{j}] $$
$$ \vec{v}_2 = 1.4 [-(0.238)\hat{i} + (0.971)\hat{j}] $$
$$ \vec{v}_2 = -0.33\hat{i} + 1.36\hat{j} \, \si{m/s} $$

\subsubsection*{4. Cálculo de Aceleración Media}
$$ \vec{a}_m = \frac{(-0.33\hat{i} + 1.36\hat{j}) - (0.98\hat{i} - 0.99\hat{j})}{3 - 1.7} $$
$$ \vec{a}_m = \frac{(-0.33 - 0.98)\hat{i} + (1.36 - (-0.99))\hat{j}}{1.3} $$
$$ \vec{a}_m = \frac{-1.31\hat{i} + 2.35\hat{j}}{1.3} $$
$$ \vec{a}_m \approx -1.0\hat{i} + 1.8\hat{j} \, \si{m/s^2} $$

\begin{tcolorbox}[colback=blue!5!white, colframe=blue!75!black, title=Respuesta Final]
\centering $\vec{a}_m = -1.0\hat{i} + 1.8\hat{j} \, \si{m/s^2}$
\end{tcolorbox}

\subsubsection*{5. Visualización Vectorial (Python)}
El siguiente código grafica las posiciones en la circunferencia y sus respectivos vectores de velocidad.

\begin{lstlisting}[language=Python]
import matplotlib.pyplot as plt
import numpy as np

# Datos
R = 0.7
omega = 2
theta_0 = np.pi/6
t1 = 1.7
t2 = 3.0

# Funciones para posicion y velocidad
def get_state(t):
    theta = theta_0 + omega * t
    # Posicion (x, y)
    rx = R * np.cos(theta)
    ry = R * np.sin(theta)
    # Velocidad (vector tangente: -sin, cos)
    v = omega * R
    vx = v * -np.sin(theta)
    vy = v * np.cos(theta)
    return rx, ry, vx, vy

# Calculo de estados
rx1, ry1, vx1, vy1 = get_state(t1)
rx2, ry2, vx2, vy2 = get_state(t2)

# Grafico
fig, ax = plt.subplots(figsize=(6, 6))

# Circulo trayectoria
circle = plt.Circle((0, 0), R, color='gray', fill=False, linestyle='--')
ax.add_artist(circle)

# Vectores Velocidad
# Escala para visualizar mejor los vectores (scale=1 dibuja magnitud real)
ax.quiver(rx1, ry1, vx1, vy1, angles='xy', scale_units='xy', scale=1, 
          color='blue', label=f'v1 (t={t1}s)')
ax.quiver(rx2, ry2, vx2, vy2, angles='xy', scale_units='xy', scale=1, 
          color='red', label=f'v2 (t={t2}s)')

# Puntos de posicion
ax.scatter([rx1, rx2], [ry1, ry2], color='black', zorder=5)
ax.text(rx1, ry1, ' P1', fontsize=10)
ax.text(rx2, ry2, ' P2', fontsize=10)

ax.set_xlim(-1.5, 1.5)
ax.set_ylim(-1.5, 1.5)
ax.set_aspect('equal')
ax.grid(True)
ax.legend()
plt.title('Vectores Velocidad en la Trayectoria')
plt.xlabel('x (m)')
plt.ylabel('y (m)')
plt.show()
\end{lstlisting}

\newpage

\subsection*{Enunciado}
Una partícula gira en el sentido del avance de las manecillas del reloj en un plano vertical $yz$ sobre una trayectoria circular de radio $R = 10 \, \si{m}$ con Movimiento Circular Uniformemente Variado (MCUV).

En el instante $t = 5 \, \si{s}$ se encuentra en el punto $\vec{r} = -10\hat{k} \, \si{m}$ y tiene una aceleración de $\vec{a} = (5\hat{k} - 4\hat{j}) \, \si{m/s^2}$.

Para el instante $t = 0$, determine los vectores aceleración angular ($\vec{\alpha}$) y velocidad angular ($\vec{\omega}_0$).

\begin{center}
\begin{tikzpicture}[scale=0.4, >=Latex]
    % Ejes (Vista Plano YZ)
    \draw[->, gray] (-12,0) -- (12,0) node[right] {$z (+)$};
    \draw[->, gray] (0,-12) -- (0,12) node[above] {$y (+)$};
    
    % Círculo
    \draw[thick] (0,0) circle (10);
    \node at (0,0) [above left] {$O$};
    
    % Punto a t=5s (-10k) -> en el eje Z negativo (izquierda en papel)
    \coordinate (P) at (-10,0);
    \fill[black] (P) circle (10pt);
    \node at (-10, 1) {$t=5s$};
    
    % Vectores Aceleración (Descomposición)
    % a = -4j + 5k
    % 5k apunta hacia el centro (positiva en z) -> a_N
    % -4j apunta hacia abajo (negativa en y) -> a_T
    
    \draw[->, thick, blue] (P) -- (-5,0) node[midway, above] {$\vec{a}_N$};
    \draw[->, thick, red] (P) -- (-10,-4) node[midway, left] {$\vec{a}_T$};
    
    % Vector total a (dashed)
    \draw[dashed, gray] (-10,-4) -- (-5,-4) -- (-5,0);
    \draw[->, purple, thick] (P) -- (-5,-4) node[below right] {$\vec{a}$};

    % Sentido de giro (Clockwise visto desde eje X positivo?)
    % La imagen sugiere a_T hacia abajo en la izquierda, esto implica que baja.
\end{tikzpicture}
\end{center}

\subsection*{Solución}

\subsubsection*{1. Identificación de Componentes}
Analizamos el vector aceleración $\vec{a} = -4\hat{j} + 5\hat{k}$ en la posición $\vec{r} = -10\hat{k}$.

\begin{itemize}
    \item \textbf{Aceleración Normal ($\vec{a}_N$):} Debe apuntar siempre hacia el centro de la trayectoria. Dado que la partícula está en $z = -10$ (eje negativo), el centro está en la dirección positiva de $z$. Por lo tanto, la componente en $\hat{k}$ corresponde a la normal:
    $$ \vec{a}_N = 5\hat{k} \, \si{m/s^2} \quad \Rightarrow \quad a_N = 5 \, \si{m/s^2} $$
    
    \item \textbf{Aceleración Tangencial ($\vec{a}_T$):} Es la componente perpendicular al radio (tangente a la trayectoria). Corresponde a la componente en $\hat{j}$:
    $$ \vec{a}_T = -4\hat{j} \, \si{m/s^2} \quad \Rightarrow \quad a_T = 4 \, \si{m/s^2} $$
\end{itemize}

\subsubsection*{2. Cálculo de Magnitudes}
Usamos las relaciones del MCUV para hallar las magnitudes angulares en $t=5\si{s}$.

$$ \omega_5 = \sqrt{\frac{a_N}{R}} = \sqrt{\frac{5}{10}} = \sqrt{0.5} \approx 0.707 \, \si{rad/s} $$

$$ \alpha = \frac{a_T}{R} = \frac{4}{10} = 0.4 \, \si{rad/s^2} $$

\subsubsection*{3. Análisis Vectorial (Regla de la Mano Derecha)}
El movimiento ocurre en el plano $yz$. Por lo tanto, los vectores angulares ($\vec{\omega}, \vec{\alpha}$) deben estar sobre el eje perpendicular, es decir, el eje $x$ (vector unitario $\hat{i}$).

\begin{itemize}
    \item \textbf{Dirección de $\vec{\alpha}$:} La aceleración tangencial apunta hacia $-y$ (abajo) en la posición $-z$ (izquierda). Esto indica que la rotación se está acelerando en ese sentido. Según la regla de la mano derecha, para generar un giro que empuje hacia $-y$ estando en $-z$, el pulgar apunta hacia adentro del plano (dirección $-\hat{i}$).
    $$ \vec{\alpha} = -0.4\hat{i} \, \si{rad/s^2} $$
    
    \item \textbf{Dirección de $\vec{\omega}_5$:} Dado que el movimiento es acelerado (la aceleración tangencial favorece el movimiento), la velocidad angular tiene el mismo signo que la aceleración angular.
    $$ \vec{\omega}_5 = -0.707\hat{i} \, \si{rad/s} $$
\end{itemize}

\subsubsection*{4. Cálculo para $t=0$}
Usamos la ecuación vectorial de la velocidad angular:
$$ \vec{\omega}_f = \vec{\omega}_0 + \vec{\alpha}(t_f - t_0) $$
$$ \vec{\omega}_0 = \vec{\omega}_5 - \vec{\alpha}(t_5 - t_0) $$

Sustituimos los vectores (trabajando solo con la componente $\hat{i}$):
$$ \vec{\omega}_0 = [-0.707\hat{i}] - [-0.4\hat{i}](5 - 0) $$
$$ \vec{\omega}_0 = -0.707\hat{i} - (-2\hat{i}) $$
$$ \vec{\omega}_0 = (-0.707 + 2)\hat{i} $$
$$ \vec{\omega}_0 = 1.293\hat{i} \, \si{rad/s} $$

\begin{tcolorbox}[colback=blue!5!white, colframe=blue!75!black, title=Respuesta Final]
\centering 
$\vec{\alpha} = -0.4\hat{i} \, \si{rad/s^2}$ \\
$\vec{\omega}_0 = 1.29\hat{i} \, \si{rad/s}$
\end{tcolorbox}

\subsubsection*{5. Visualización 3D (Python)}
Usamos una gráfica 3D para visualizar que los vectores angulares son perpendiculares al plano de rotación.

\begin{lstlisting}[language=Python]
import matplotlib.pyplot as plt
import numpy as np

# Datos
R = 10
alpha_vec = np.array([-0.4, 0, 0])
omega0_vec = np.array([1.29, 0, 0])

fig = plt.figure(figsize=(8, 6))
ax = fig.add_subplot(111, projection='3d')

# Dibujar circulo en plano YZ
theta = np.linspace(0, 2*np.pi, 100)
y = R * np.cos(theta)
z = R * np.sin(theta)
x = np.zeros_like(theta)
ax.plot(x, z, y, color='gray', linestyle='--', alpha=0.5, label='Trayectoria (Plano YZ)')

# Ejes
ax.quiver(0, 0, 0, 15, 0, 0, color='k', arrow_length_ratio=0.1) # X
ax.quiver(0, 0, 0, 0, 15, 0, color='k', arrow_length_ratio=0.1) # Z (ajuste visual matplotlib)
ax.quiver(0, 0, 0, 0, 0, 15, color='k', arrow_length_ratio=0.1) # Y
ax.text(16, 0, 0, 'x')
ax.text(0, 16, 0, 'z')
ax.text(0, 0, 16, 'y')

# Vectores Angulares (En el eje X)
# Omega 0
ax.quiver(0, 0, 0, omega0_vec[0]*5, 0, 0, color='green', linewidth=2, 
          label=r'$\vec{\omega}_0$ (Escalado)')
# Alpha
ax.quiver(0, 0, 0, alpha_vec[0]*5, 0, 0, color='red', linewidth=2, 
          label=r'$\vec{\alpha}$ (Escalado)')

# Punto en t=5 (-10k) -> z=-10, y=0
ax.scatter(0, -10, 0, color='blue', s=50, label='Posicion t=5s')
# Vector aceleracion en t=5 (a = -4j + 5k)
# En plot 3d: (x, z, y) -> (0, -10, 0) origen del vector
# Vector componentes: dx=0, dz=5, dy=-4
ax.quiver(0, -10, 0, 0, 5, -4, color='purple', length=1.5, label=r'$\vec{a}(t=5)$')

ax.set_xlabel('X')
ax.set_ylabel('Z')
ax.set_zlabel('Y')
ax.set_title('Cinematica Rotacional en 3D')
ax.legend()
plt.show()
\end{lstlisting}

\newpage

\subsection*{Enunciado}
La posición angular de una partícula que describe una trayectoria circular de radio $R = 1.2 \, \si{m}$, en sentido contrario al avance de las manecillas del reloj, viene dada por la función:
$$ \theta(t) = \frac{\pi}{2} + t + t^2 $$
Donde $\theta$ está en radianes y $t$ en segundos. Determine la ecuación de la velocidad angular $\omega$ en función del tiempo.

\begin{center}
\begin{tikzpicture}[scale=1.5, >=Latex]
    % Ejes
    \draw[->, gray] (-1.5,0) -- (1.5,0) node[right] {$x$};
    \draw[->, gray] (0,-1.5) -- (0,1.5) node[above] {$y$};
    
    % Trayectoria
    \draw[thick] (0,0) circle (1.2);
    
    % Radio vector genérico
    \draw[->, thick, black] (0,0) -- (60:1.2) node[midway, above left] {$R=1.2$};
    
    % Dirección de movimiento (Antihorario)
    \draw[->, blue, thick] (1.4, 0.2) arc (0:60:1.4) node[midway, right] {$\omega(+)$};
    
    % Etiqueta de la función
    \node[blue, right] at (0.5, -1.6) {$\theta(t) = \frac{\pi}{2} + t + t^2$};
\end{tikzpicture}
\end{center}

\subsection*{Solución}

\subsubsection*{1. Concepto: Derivada de la Posición}
La velocidad angular instantánea $\omega(t)$ se define como la derivada de la posición angular $\theta(t)$ con respecto al tiempo:
$$ \omega(t) = \frac{d\theta}{dt} $$

\subsubsection*{2. Derivación}
Dada la función:
$$ \theta(t) = \frac{\pi}{2} + t + t^2 $$

Derivamos término a término:
\begin{itemize}
    \item La derivada de una constante ($\pi/2$) es $0$.
    \item La derivada de $t$ es $1$.
    \item La derivada de $t^2$ es $2t$.
\end{itemize}

Por lo tanto:
$$ \omega(t) = 0 + 1 + 2t $$
$$ \boxed{\omega(t) = 1 + 2t \, (\si{rad/s})} $$

\subsubsection*{3. Identificación de Parámetros}
Comparando con la ecuación general del Movimiento Circular Uniformemente Variado (MCUV):
$$ \omega(t) = \omega_0 + \alpha t $$
Podemos identificar que:
\begin{itemize}
    \item Velocidad angular inicial: $\omega_0 = 1 \, \si{rad/s}$
    \item Aceleración angular: $\alpha = 2 \, \si{rad/s^2}$
\end{itemize}

\subsubsection*{4. Gráfica de la Función (Python)}
\begin{lstlisting}[language=Python]
import matplotlib.pyplot as plt
import numpy as np

# Tiempo
t = np.linspace(0, 5, 100)

# Funciones cinemáticas
theta = np.pi/2 + t + t**2
omega = 1 + 2*t

fig, (ax1, ax2) = plt.subplots(1, 2, figsize=(10, 4))

# Grafica Posicion vs Tiempo
ax1.plot(t, theta, 'purple', label=r'$\theta(t) = \pi/2 + t + t^2$')
ax1.set_xlabel('Tiempo (s)')
ax1.set_ylabel('Posicion Angular (rad)')
ax1.set_title('Posicion Angular')
ax1.grid(True)
ax1.legend()

# Grafica Velocidad vs Tiempo
ax2.plot(t, omega, 'blue', label=r'$\omega(t) = 1 + 2t$')
ax2.set_xlabel('Tiempo (s)')
ax2.set_ylabel('Velocidad Angular (rad/s)')
ax2.set_title('Velocidad Angular')
ax2.grid(True)
ax2.legend()

plt.tight_layout()
plt.show()
\end{lstlisting}

\newpage

\subsection*{Enunciado}
La posición angular de una partícula que describe una trayectoria circular de radio $R = 1.2 \, \si{m}$, en sentido contrario al avance de las manecillas del reloj, viene dada por la función:
$$ \theta(t) = \frac{\pi}{2} + t + t^2 $$
Donde $\theta$ está en radianes y $t$ en segundos. Determine el vector aceleración total $\vec{a}$ en el instante $t = 3 \, \si{s}$.

\begin{center}
\begin{tikzpicture}[scale=1.5, >=Latex]
    % Ejes
    \draw[->, gray] (-1.5,0) -- (1.5,0) node[right] {$x$};
    \draw[->, gray] (0,-1.5) -- (0,1.5) node[above] {$y$};
    \draw[thick] (0,0) circle (1.2);
    
    % Representación conceptual (no a escala exacta para t=3, solo esquema)
    \coordinate (P) at (60:1.2);
    \fill (P) circle (2pt) node[above right] {Partícula en $t=3s$};
    \draw[->, thick] (0,0) -- (P);
    
    % Vectores aceleración componentes
    \draw[->, red, thick] (P) -- ++(200:0.8) node[midway, below] {$\vec{a}_N$};
    \draw[->, blue, thick] (P) -- ++(150:0.8) node[midway, left] {$\vec{a}_T$};
\end{tikzpicture}
\end{center}

\subsection*{Solución}

\subsubsection*{1. Estado Cinemático en $t=3\,\si{s}$}
Primero calculamos la posición angular y la velocidad angular en el instante solicitado.

\textbf{Posición angular ($\theta_3$):}
$$ \theta(3) = \frac{\pi}{2} + 3 + (3)^2 = 1.57 + 3 + 9 = 13.57 \, \si{rad} $$
Para ubicar el vector, convertimos a grados y reducimos al primer círculo ($0^\circ - 360^\circ$):
$$ 13.57 \, \si{rad} \times \frac{180}{\pi} \approx 777.5^\circ $$
$$ 777.5^\circ - 2(360^\circ) = 57.5^\circ $$
El ángulo efectivo respecto al eje X positivo es $\phi \approx 57.9^\circ$ (usando el valor de referencia del texto).

\textbf{Velocidad angular ($\omega_3$):}
$$ \omega(3) = 1 + 2(3) = 7 \, \si{rad/s} $$

\textbf{Aceleración angular ($\alpha$):}
$$ \alpha = 2 \, \si{rad/s^2} \quad (\text{Constante}) $$

\subsubsection*{2. Magnitudes de Aceleración}
Calculamos las componentes intrínsecas (escalares):

\begin{itemize}
    \item \textbf{Aceleración Tangencial ($a_T$):}
    $$ a_T = \alpha R = (2)(1.2) = 2.4 \, \si{m/s^2} $$
    
    \item \textbf{Aceleración Normal ($a_N$):}
    $$ a_N = \omega^2 R = (7)^2 (1.2) = 49(1.2) = 58.8 \, \si{m/s^2} $$
\end{itemize}

\subsubsection*{3. Descomposición Vectorial}
La partícula está en el primer cuadrante ($57.9^\circ$).
\begin{itemize}
    \item El vector unitario radial es $\hat{u}_r = \cos\phi \, \hat{i} + \sin\phi \, \hat{j}$.
    \item El vector unitario tangencial (sentido antihorario) es $\hat{u}_\theta = -\sin\phi \, \hat{i} + \cos\phi \, \hat{j}$.
\end{itemize}

\textbf{Vector Aceleración Tangencial ($\vec{a}_T$):}
Va en la dirección del movimiento (tangente).
$$ \vec{a}_T = 2.4 (-\sin 57.9^\circ \hat{i} + \cos 57.9^\circ \hat{j}) $$
$$ \vec{a}_T = 2.4 (-0.847 \hat{i} + 0.531 \hat{j}) $$
$$ \vec{a}_T \approx -2.03\hat{i} + 1.27\hat{j} \, \si{m/s^2} $$

\textbf{Vector Aceleración Normal ($\vec{a}_N$):}
Va hacia el centro (opuesto al radial).
$$ \vec{a}_N = 58.8 (-\cos 57.9^\circ \hat{i} - \sin 57.9^\circ \hat{j}) $$
$$ \vec{a}_N = 58.8 (-0.531 \hat{i} - 0.847 \hat{j}) $$
$$ \vec{a}_N \approx -31.22\hat{i} - 49.80\hat{j} \, \si{m/s^2} $$

\subsubsection*{4. Vector Aceleración Total ($\vec{a}$)}
Sumamos vectorialmente $\vec{a} = \vec{a}_T + \vec{a}_N$:
$$ \vec{a} = (-2.03 - 31.22)\hat{i} + (1.27 - 49.80)\hat{j} $$
$$ \vec{a} = -33.25\hat{i} - 48.53\hat{j} \, \si{m/s^2} $$

\begin{tcolorbox}[colback=blue!5!white, colframe=blue!75!black, title=Respuesta Final]
\centering $\vec{a} \approx -33.2\hat{i} - 48.5\hat{j} \, \si{m/s^2}$
\end{tcolorbox}

\newpage

\newpage

\subsection*{Enunciado}
Una partícula se mueve por una trayectoria circular de radio $R = 2 \, \si{m}$ con movimiento uniformemente variado. En el instante $t = 0$, su velocidad es $\vec{v}_0 = 3\hat{i} + 4\hat{j} \, \si{m/s}$ y su rapidez disminuye a razón de $2 \, \si{m/s}$ cada segundo.
Determine el vector posición de la partícula en $t=0$.

\begin{center}
\begin{tikzpicture}[scale=1.5, >=Latex]
    % Ejes
    \draw[->, gray] (-2.5,0) -- (2.5,0) node[right] {$x$};
    \draw[->, gray] (0,-2.5) -- (0,2.5) node[above] {$y$};
    
    % Círculo
    \draw[thick] (0,0) circle (2);
    
    % Vector Velocidad (esquemático en 1er cuadrante)
    % 3i + 4j es ángulo 53 grados. Posición debe ser 53-90 = -37 grados
    \coordinate (P) at (-36.87:2);
    \draw[dashed] (0,0) -- (P) node[midway, below left] {$R=2$};
    \fill (P) circle (2pt) node[right] {$A (t=0)$};
    
    % Vector v0 (tangente)
    \draw[->, blue, thick] (P) -- ++(53.13:1.5) node[above] {$\vec{v}_0$};
    
    % Flecha de movimiento (antihorario)
    \draw[->] (2.2, -0.5) arc (-10:10:2.2) node[midway, right] {mov};
\end{tikzpicture}
\end{center}

\subsection*{Solución}

\subsubsection*{1. Análisis Geométrico}
Sabemos que en un movimiento circular, el vector posición $\vec{r}$ es perpendicular al vector velocidad $\vec{v}$.
Dado que $\vec{v}_0 = 3\hat{i} + 4\hat{j}$, calculamos su ángulo de inclinación $\beta$:
$$ \beta = \tan^{-1}\left(\frac{4}{3}\right) \approx 53.13^\circ $$

Como el movimiento es antihorario (levógiro), el vector posición se encuentra $90^\circ$ "atrás" del vector velocidad (en sentido horario respecto a la velocidad).
$$ \theta_{pos} = 53.13^\circ - 90^\circ = -36.87^\circ $$

\subsubsection*{2. Cálculo Vectorial}
El vector posición tiene magnitud $R=2$ y dirección $\theta_{pos}$.
\begin{align*}
    \vec{r}_0 &= R (\cos \theta_{pos} \hat{i} + \sin \theta_{pos} \hat{j}) \\
    \vec{r}_0 &= 2 (\cos(-36.87^\circ)\hat{i} + \sin(-36.87^\circ)\hat{j})
\end{align*}
Usando las propiedades trigonométricas ($\cos(-x)=\cos x$, $\sin(-x)=-\sin x$) y la relación con el ángulo $\beta$:
$$ \vec{r}_0 = 2 (\sin 53.13^\circ \hat{i} - \cos 53.13^\circ \hat{j}) $$
Sabemos que $\sin 53.13^\circ = 0.8$ y $\cos 53.13^\circ = 0.6$:
$$ \vec{r}_0 = 2 (0.8 \hat{i} - 0.6 \hat{j}) $$
$$ \vec{r}_0 = 1.6\hat{i} - 1.2\hat{j} \, \si{m} $$

\begin{tcolorbox}[colback=blue!5!white, colframe=blue!75!black, title=Respuesta Final]
\centering $\vec{r}_0 = 1.6\hat{i} - 1.2\hat{j} \, \si{m}$
\end{tcolorbox}

\subsubsection*{3. Visualización (Python)}
\begin{lstlisting}[language=Python]
import matplotlib.pyplot as plt
import numpy as np

# Datos
v0 = np.array([3, 4])
R = 2

# Calculo de posicion (perpendicular a v0, sentido horario relativo)
# Rotacion de -90 grados (matriz de rotacion)
# x' = x*cos(-90) - y*sin(-90) = y
# y' = x*sin(-90) + y*cos(-90) = -x
# direccion unitaria de r: (vy, -vx) normalizado? 
# No, para giro antihorario, r esta a la derecha de v.
# v=(3,4). r debe ser proporcional a (4, -3) o (-4, 3)?
# Si r está en 4to cuadrante (1.6, -1.2), r es proporcional a (4, -3).

r0_x = 1.6
r0_y = -1.2

fig, ax = plt.subplots(figsize=(5, 5))
circle = plt.Circle((0, 0), R, color='gray', fill=False, linestyle='--')
ax.add_artist(circle)

# Vectores
ax.quiver(0, 0, r0_x, r0_y, angles='xy', scale_units='xy', scale=1, 
          color='green', label=r'$\vec{r}_0$')
ax.quiver(r0_x, r0_y, v0[0], v0[1], angles='xy', scale_units='xy', scale=1, 
          color='blue', label=r'$\vec{v}_0$')

# Punto A
ax.scatter(r0_x, r0_y, color='black', zorder=5)
ax.text(r0_x + 0.2, r0_y, 'A (t=0)')

ax.set_xlim(-3, 4)
ax.set_ylim(-3, 5)
ax.set_aspect('equal')
ax.grid(True)
plt.legend()
plt.title('Posición y Velocidad Inicial')
plt.show()
\end{lstlisting}

\newpage

\subsection*{Enunciado}
Una partícula describe un MCUV con radio $R=2\,\si{m}$. En $t=0$, su rapidez es $v_0 = 5\,\si{m/s}$ y frena a razón de $2\,\si{m/s}$ cada segundo.
Determine el vector aceleración total en $t = 2.5\,\si{s}$.

\subsection*{Solución}

\subsubsection*{1. Estado del Movimiento en $t=2.5s$}
Calculamos la rapidez lineal en ese instante:
$$ v(t) = v_0 + a_t t $$
Dado que frena, $a_t = -2\,\si{m/s^2}$.
$$ v(2.5) = 5 + (-2)(2.5) = 5 - 5 = 0 \, \si{m/s} $$

La partícula se detiene momentáneamente. Esto implica que la velocidad angular $\omega = 0$.

\subsubsection*{2. Componentes de la Aceleración}
\begin{itemize}
    \item \textbf{Aceleración Normal:} Depende de la velocidad al cuadrado.
    $$ a_N = \frac{v^2}{R} = \frac{0^2}{2} = 0 \, \si{m/s^2} $$
    \item \textbf{Aceleración Tangencial:} Es constante dada por el problema.
    $$ a_T = 2 \, \si{m/s^2} \quad (\text{magnitud}) $$
\end{itemize}
La aceleración total es igual a la aceleración tangencial en este instante: $\vec{a}_{2.5} = \vec{a}_T$.

\subsubsection*{3. Posición en $t=2.5s$}
Necesitamos la posición para saber la dirección del vector tangencial.
$$ \theta(t) = \theta_0 + \omega_0 t + \frac{1}{2}\alpha t^2 $$
Datos iniciales:
\begin{itemize}
    \item $\theta_0 = -36.87^\circ \approx -0.64 \, \si{rad}$ (del Ejercicio 6).
    \item $\omega_0 = v_0/R = 5/2 = 2.5 \, \si{rad/s}$.
    \item $\alpha = a_t/R = -2/2 = -1 \, \si{rad/s^2}$.
\end{itemize}

$$ \theta(2.5) = -0.64 + (2.5)(2.5) + \frac{1}{2}(-1)(2.5)^2 $$
$$ \theta(2.5) = -0.64 + 6.25 - 3.125 = 2.485 \, \si{rad} $$
Convertimos a grados: $2.485 \times (180/\pi) \approx 142.4^\circ$. (El texto original usa otra referencia angular, pero físicamente $142.4^\circ$ es la posición correcta en el 2do cuadrante).

\subsubsection*{4. Vector Aceleración}
La aceleración tangencial se opone al movimiento original (antihorario), por lo que apunta en sentido horario (tangente hacia "atrás").
Vector unitario tangencial en $\theta = 142.4^\circ$:
$$ \hat{u}_T = -\sin(142.4^\circ)\hat{i} + \cos(142.4^\circ)\hat{j} $$
La aceleración tangencial $\vec{a}_T$ es el vector que causa el frenado. Si asumimos la dirección positiva antihoraria, $a_t$ es negativo, apuntando a favor de $-\hat{u}_T$ (horario).
$$ \vec{a} = 2 (\cos(142.4^\circ - 90^\circ)\hat{i} + \sin(142.4^\circ - 90^\circ)\hat{j}) $$
$$ \vec{a} = 2 (\cos 52.4^\circ \hat{i} + \sin 52.4^\circ \hat{j}) $$
$$ \vec{a} \approx 1.22\hat{i} + 1.58\hat{j} \, \si{m/s^2} $$
*(Nota: Valores aproximados coinciden con $1.24\hat{i} + 1.56\hat{j}$ del texto original usando redondeos intermedios).*

\begin{tcolorbox}[colback=blue!5!white, colframe=blue!75!black, title=Respuesta Final]
\centering $\vec{a} \approx 1.24\hat{i} + 1.56\hat{j} \, \si{m/s^2}$
\end{tcolorbox}

\newpage

\subsection*{Enunciado}
Para la misma partícula ($v_0 = 3\hat{i} + 4\hat{j}$ en $t=0$, $a_t = -2\,\si{m/s^2}$), determine el vector velocidad en $t=5\,\si{s}$.

\subsection*{Solución}

\subsubsection*{1. Simetría del Movimiento}
El tiempo de subida (hasta detenerse) fue $t_{stop} = 2.5\,\si{s}$.
El tiempo $t=5\,\si{s}$ es exactamente el doble ($2 \times t_{stop}$).
Por simetría del movimiento uniformemente variado, la partícula regresa a la posición angular inicial.

\subsubsection*{2. Verificación de Posición}
$$ \Delta \theta = \omega_0 t + \frac{1}{2}\alpha t^2 $$
$$ \Delta \theta = (2.5)(5) + \frac{1}{2}(-1)(25) $$
$$ \Delta \theta = 12.5 - 12.5 = 0 \, \si{rad} $$
La partícula está en la misma posición angular que en $t=0$.

\subsubsection*{3. Vector Velocidad}
Calculamos la rapidez escalar:
$$ v(5) = 5 + (-2)(5) = -5 \, \si{m/s} $$
La magnitud es la misma ($5\,\si{m/s}$), pero el signo negativo indica que se mueve en sentido opuesto (horario).

Si en $t=0$ el vector era $\vec{v}_0$, en el mismo punto pero sentido opuesto, el vector es:
$$ \vec{v}_5 = -\vec{v}_0 $$
$$ \vec{v}_5 = -(3\hat{i} + 4\hat{j}) $$

\begin{tcolorbox}[colback=green!5!white, colframe=green!50!black, title=Respuesta Final]
\centering $\vec{v}_5 = -3\hat{i} - 4\hat{j} \, \si{m/s}$
\end{tcolorbox}

\newpage

\subsection*{Enunciado}
Una partícula describe un MCUV con radio $R=2\,\si{m}$. En $t=0$, su rapidez es $v_0 = 5\,\si{m/s}$ y frena a razón de $2\,\si{m/s}$ cada segundo. Determine el vector aceleración tangencial $\vec{a}_T$ para los instantes $t=0$ y $t=5\,\si{s}$.

\subsection*{Solución}

\subsubsection*{1. Concepto Físico}
La aceleración tangencial es responsable del cambio en la magnitud de la velocidad. Su magnitud es constante $|\vec{a}_T| = 2\,\si{m/s^2}$.
Su dirección depende de la posición en la trayectoria.

\subsubsection*{2. Instante $t=0$}
La partícula está en $A$ (4to cuadrante). Se mueve antihorario, pero frena. Por tanto, $\vec{a}_T$ apunta en sentido horario (opuesto a $\vec{v}_0$).
$$ \vec{v}_0 = 3\hat{i} + 4\hat{j} $$
El vector opuesto unitario es:
$$ \hat{u}_{aT} = \frac{-3\hat{i} - 4\hat{j}}{5} = -0.6\hat{i} - 0.8\hat{j} $$
Multiplicamos por la magnitud (2):
$$ \vec{a}_T(0) = 2(-0.6\hat{i} - 0.8\hat{j}) = -1.2\hat{i} - 1.6\hat{j} \, \si{m/s^2} $$

\subsubsection*{3. Instante $t=5$}
La partícula está nuevamente en $A$. Se mueve en sentido horario (velocidad negativa) y su rapidez "negativa" sigue aumentando en magnitud (o físicamente, el vector aceleración angular sigue siendo el mismo $\vec{\alpha} = -1\hat{k}$).
La fuerza tangencial no ha cambiado de dirección en el espacio; sigue "empujando" en sentido horario.

Por tanto, el vector aceleración tangencial es el mismo que en $t=0$:
$$ \vec{a}_T(5) = -1.2\hat{i} - 1.6\hat{j} \, \si{m/s^2} $$

\begin{tcolorbox}[colback=blue!5!white, colframe=blue!75!black, title=Respuesta Final]
\centering $\vec{a}_T = -1.2\hat{i} - 1.6\hat{j} \, \si{m/s^2}$ (para ambos casos)
\end{tcolorbox}

\subsubsection*{4. Gráfico de Aceleración (Python)}
\begin{lstlisting}[language=Python]
import matplotlib.pyplot as plt

# Vector aceleracion tangencial
at_x = -1.2
at_y = -1.6
r0_x = 1.6
r0_y = -1.2

plt.figure(figsize=(5, 5))

# Posicion A
plt.scatter(r0_x, r0_y, color='black', label='Posición A')
plt.plot([0, r0_x], [0, r0_y], 'k--', alpha=0.3)

# Vector v0 (t=0)
plt.quiver(r0_x, r0_y, 3, 4, angles='xy', scale_units='xy', scale=1, 
           color='blue', alpha=0.3, label='v0 (t=0)')

# Vector a_T (t=0 y t=5)
# Nota: Se dibuja desde la posicion de la particula
plt.quiver(r0_x, r0_y, at_x, at_y, angles='xy', scale_units='xy', scale=1, 
           color='red', label=r'$\vec{a}_T$')

plt.xlim(-2, 5)
plt.ylim(-4, 4)
plt.grid(True)
plt.legend()
plt.title('Vector Aceleración Tangencial en A')
plt.gca().set_aspect('equal')
plt.show()
\end{lstlisting}

\newpage

\subsection*{Enunciado}
Las ruedas $A$ y $B$, de $2$ y $6 \, \si{cm}$ de radio, están interconectadas como indica la figura. Si la velocidad angular de $A$ es $4\vec{k} \, \si{rad/s}$, determine la aceleración centrípeta de la partícula $P$ en el instante mostrado.

\begin{center}
\begin{tikzpicture}[scale=1, >=Latex]
    \draw[thick] (0,0) circle (0.6);
    \draw[thick] (2.4,0) circle (1.8);
    
    \node at (0.3,0.3) {$A$};
    \node at (3.5,1) {$B$};
    
    \draw[dashed] (0,-0.6) -- (0,0.6);
    \draw[dashed] (-0.6,0) -- (0.6,0);
    
    \draw[dashed] (2.4,-1.8) -- (2.4,1.8);
    \draw[dashed] (0.6,0) -- (4.2,0);
    
    \node[below] at (2.4,-1.8) {$P$};
    
    \draw[->, thick] (-0.5,0.7) arc (135:45:0.6);
\end{tikzpicture}
\end{center}

\subsection*{Solución}

\subsubsection*{1. Análisis de la transmisión de movimiento}
Dado que las ruedas $A$ y $B$ están interconectadas por contacto directo en sus bordes sin deslizar, la velocidad tangencial en el punto de contacto es la misma para ambas ruedas.
$$ v_A = v_B $$

\subsubsection*{2. Cálculo de la velocidad tangencial}
Primero, convertimos el radio de la rueda $A$ al Sistema Internacional (SI):
$$ r_A = 2 \, \si{cm} = 0.02 \, \si{m} $$

La magnitud de la velocidad angular de $A$ es $\omega_A = 4 \, \si{rad/s}$. Calculamos su velocidad tangencial:
$$ v_A = \omega_A r_A $$
$$ v_A = (4)(0.02) $$
$$ v_A = 0.08 \, \si{m/s} $$

Por la relación de transmisión establecida en el paso anterior, la velocidad tangencial en el borde de la rueda $B$ también es:
$$ v_B = 0.08 \, \si{m/s} $$

\subsubsection*{3. Análisis de la partícula P}
La partícula $P$ se encuentra en el borde exterior de la rueda $B$, por lo que su velocidad tangencial es $v_P = v_B = 0.08 \, \si{m/s}$. 

Convertimos el radio de la rueda $B$ al SI:
$$ r_B = 6 \, \si{cm} = 0.06 \, \si{m} $$

Para entender la dirección de la aceleración centrípeta, graficamos el vector en la rueda $B$:

\begin{lstlisting}
import matplotlib.pyplot as plt
import matplotlib.patches as patches

fig, ax = plt.subplots(figsize=(4, 4))
ax.set_xlim(-2, 2); ax.set_ylim(-2, 2)
ax.axis('off')

circleB = patches.Circle((0, 0), 1.5, fill=False, edgecolor='black', linewidth=1.5)
ax.add_patch(circleB)
ax.text(0.3, 0.3, 'B', fontsize=12)

ax.plot(0, -1.5, 'ko')
ax.text(0, -1.7, 'P', fontsize=12, ha='center')

ax.arrow(0, -1.5, 0, 0.8, head_width=0.1, fc='blue', ec='blue')
ax.text(0.1, -0.9, r'$\vec{a}_c$', color='blue', fontsize=12)

plt.tight_layout()
plt.show()
\end{lstlisting}

La aceleración centrípeta ($a_c$) siempre apunta hacia el centro de rotación. Como la partícula $P$ está en la parte inferior de la rueda $B$, el vector apuntará verticalmente hacia arriba, es decir, en la dirección del vector unitario $\vec{j}$.

Calculamos su magnitud:
$$ a_c = \frac{v_P^2}{r_B} $$
$$ a_c = \frac{(0.08)^2}{0.06} $$
$$ a_c = \frac{0.0064}{0.06} $$
$$ a_c \approx 0.1067 \, \si{m/s^2} $$

Expresando el resultado en forma vectorial:
$$ \vec{a}_c = 0.1067\vec{j} \, \si{m/s^2} $$

\begin{tcolorbox}[colback=green!5!white, colframe=green!50!black, title=Respuesta Final]
\centering \Large $\vec{a}_c = 0.107\vec{j} \, \si{m/s^2}$
\end{tcolorbox}

\newpage

\subsection*{Enunciado}
Un niño se encuentra sobre el borde de un carrusel de $12 \, \si{m}$ de radio (fig.), que gira a razón de $0.5\vec{k} \, \si{rad/s}$. Determine el tiempo que le toma retornar a su puesto, si el niño camina por el borde a razón de $1 \, \si{m/s}$ respecto a tierra, en sentido contrario al carrusel. 

\begin{center}
\begin{tikzpicture}[scale=1, >=Latex]
    \draw[->, thick] (-2.5,0) -- (2.5,0) node[right] {$x$};
    \draw[->, thick] (0,-2.5) -- (0,2.5) node[above] {$y$};
    
    \draw[thick] (0,0) circle (2);
    
    \fill (2,0) circle (0.08);
    \node[above left] at (2.2,0) {niño};
    
    \draw[->, thick, dashed] (2.2,0) arc (0:-60:2.2);
\end{tikzpicture}
\end{center}

\subsection*{Solución}

\subsubsection*{1. Análisis de velocidades angulares}
Se establecen los datos iniciales proporcionados por el problema:
\begin{itemize}
    \item Radio del carrusel: $R = 12 \, \si{m}$
    \item Velocidad angular del carrusel: $\vec{\omega}_C = 0.5\vec{k} \, \si{rad/s}$
    \item Rapidez del niño respecto a tierra: $v_N = 1 \, \si{m/s}$
\end{itemize}

Se calcula la magnitud de la velocidad angular del niño respecto a tierra:
$$ \omega_N = \frac{v_N}{R} $$
$$ \omega_N = \frac{1}{12} \, \si{rad/s} $$

El problema indica que el niño camina en sentido contrario al carrusel. Dado que el carrusel gira en sentido antihorario ($+\vec{k}$), el niño se mueve en sentido horario ($-\vec{k}$). Vectorialmente:
$$ \vec{\omega}_N = -\frac{1}{12}\vec{k} \, \si{rad/s} $$

\subsubsection*{2. Velocidad angular relativa}
Para determinar cuánto tarda el niño en dar una vuelta completa respecto al carrusel, se requiere calcular la velocidad angular del niño respecto al carrusel ($\vec{\omega}_{N/C}$):
$$ \vec{\omega}_{N/C} = \vec{\omega}_N - \vec{\omega}_C $$
$$ \vec{\omega}_{N/C} = \left( -\frac{1}{12}\vec{k} \right) - \left( 0.5\vec{k} \right) $$
$$ \vec{\omega}_{N/C} = \left( -\frac{1}{12} - \frac{1}{2} \right)\vec{k} $$
$$ \vec{\omega}_{N/C} = -\frac{7}{12}\vec{k} \, \si{rad/s} \approx -0.583\vec{k} \, \si{rad/s} $$

\subsubsection*{3. Cálculo del tiempo}
Para que el niño retorne a su posición inicial en el carrusel, debe describir un desplazamiento angular relativo equivalente a una vuelta completa. Como su movimiento relativo es en sentido horario, el desplazamiento es negativo:
$$ \Delta\theta = -2\pi \, \si{rad} $$

Se utiliza la ecuación de movimiento circular uniforme para calcular el tiempo:
$$ \Delta t = \frac{\Delta\theta}{\omega_{N/C}} $$
$$ \Delta t = \frac{-2\pi}{-7/12} $$
$$ \Delta t = \frac{24\pi}{7} \, \si{s} $$
$$ \Delta t \approx 10.77 \, \si{s} $$

\begin{tcolorbox}[colback=green!5!white, colframe=green!50!black, title=Respuesta Final]
\centering \Large $\Delta t = 10.77 \, \si{s}$
\end{tcolorbox}

\newpage

\subsection*{Enunciado}
En una laptop, con disco duro de $500 \, \si{GB}$ de capacidad, se tiene una rapidez angular de $5400 \, \si{rpm}$. Si la velocidad de almacenamiento de este disco duro es de $77 \, \si{MB}$ en cada segundo, determine el número de revoluciones giradas al almacenar $770 \, \si{MB}$ de información.

\subsection*{Solución}

\subsubsection*{1. Cálculo del tiempo de almacenamiento}
Se determinan los datos de almacenamiento:
\begin{itemize}
    \item Velocidad de transferencia: $v = 77 \, \si{MB/s}$
    \item Información a almacenar: $I = 770 \, \si{MB}$
\end{itemize}

Se calcula el tiempo $\Delta t$ necesario para procesar dicha información:
$$ \Delta t = \frac{I}{v} $$
$$ \Delta t = \frac{770}{77} $$
$$ \Delta t = 10 \, \si{s} $$

\subsubsection*{2. Conversión de unidades}
Se convierte el tiempo obtenido a minutos para operar con la rapidez angular dada en revoluciones por minuto:
$$ \Delta t = 10 \, \si{s} \times \left( \frac{1 \, \si{min}}{60 \, \si{s}} \right) $$
$$ \Delta t = \frac{1}{6} \, \si{min} $$

\subsubsection*{3. Cálculo del número de revoluciones}
Se utiliza la magnitud de la rapidez angular del disco duro:
$$ \omega = 5400 \, \si{rpm} = 5400 \, \frac{\si{rev}}{\si{min}} $$

Se aplica la ecuación de desplazamiento angular para movimiento circular uniforme:
$$ \Delta \theta = \omega \cdot \Delta t $$
$$ \Delta \theta = 5400 \times \left( \frac{1}{6} \right) $$
$$ \Delta \theta = \frac{5400}{6} $$
$$ \Delta \theta = 900 \, \si{rev} $$

\begin{tcolorbox}[colback=green!5!white, colframe=green!50!black, title=Respuesta Final]
\centering \Large $\Delta \theta = 900 \, \si{rev}$
\end{tcolorbox}

\newpage

\subsection*{Enunciado}
La partícula de la figura se traslada con MCUVA en sentido antihorario, por una trayectoria de $10 \, \si{m}$ de radio. El módulo de su aceleración inicial fue $4 \, \si{m/s^2}$ y su velocidad inicial $\vec{v}_0 = 3\vec{i} - 4\vec{j} \, \si{m/s}$. Para $t = 5 \, \si{s}$, determine la posición de la partícula.

\begin{center}
\begin{tikzpicture}[scale=0.8, >=Latex]
    \draw[->, thick] (-3,0) -- (3,0) node[right] {$x$};
    \draw[->, thick] (0,-3) -- (0,3) node[above] {$y$};
    
    \draw[thick] (0,0) circle (2);
    
    \fill (-1.6,-1.2) circle (0.08);
    \draw[thick, ->] (0,0) -- (-1.6,-1.2);
    
    \draw[thick, ->] (-1.6,-1.2) -- (-0.6,-2.6) node[below] {$\vec{v}_0$};
\end{tikzpicture}
\end{center}

\subsection*{Solución}

\subsubsection*{1. Análisis de las condiciones iniciales}
Se calcula el módulo de la velocidad inicial:
$$ |\vec{v}_0| = \sqrt{3^2 + (-4)^2} = \sqrt{9 + 16} = 5 \, \si{m/s} $$

Se obtiene la velocidad angular inicial ($\omega_0$):
$$ \omega_0 = \frac{v_0}{R} = \frac{5}{10} = 0.5 \, \si{rad/s} $$
Al ser el movimiento antihorario, vectorialmente es $\vec{\omega}_0 = 0.5\vec{k} \, \si{rad/s}$.

Para hallar la posición inicial ($\vec{r}_0$), se usa la relación $\vec{v}_0 = \vec{\omega}_0 \times \vec{r}_0$. Si $\vec{r}_0 = x\vec{i} + y\vec{j}$:
$$ 3\vec{i} - 4\vec{j} = (0.5\vec{k}) \times (x\vec{i} + y\vec{j}) $$
$$ 3\vec{i} - 4\vec{j} = -0.5y\vec{i} + 0.5x\vec{j} $$
Igualando componentes:
$$ -0.5y = 3 \Rightarrow y = -6 $$
$$ 0.5x = -4 \Rightarrow x = -8 $$
$$ \vec{r}_0 = -8\vec{i} - 6\vec{j} \, \si{m} $$

El ángulo de la posición inicial respecto al eje $+x$ es:
$$ \theta_0 = \arctan\left(\frac{-6}{-8}\right) $$
Como está en el tercer cuadrante:
$$ \theta_0 = 180^\circ + 36.87^\circ = 216.87^\circ = 3.785 \, \si{rad} $$

\subsubsection*{2. Cálculo de la aceleración angular}
Se calcula la aceleración centrípeta inicial ($a_c$):
$$ a_c = \frac{v_0^2}{R} = \frac{5^2}{10} = 2.5 \, \si{m/s^2} $$

El módulo de la aceleración total es $a = 4 \, \si{m/s^2}$. Se deduce la aceleración tangencial ($a_t$):
$$ a^2 = a_c^2 + a_t^2 $$
$$ 4^2 = 2.5^2 + a_t^2 $$
$$ 16 = 6.25 + a_t^2 $$
$$ a_t = \sqrt{9.75} \approx 3.122 \, \si{m/s^2} $$

La magnitud de la aceleración angular ($\alpha$) constante es:
$$ \alpha = \frac{a_t}{R} = \frac{3.122}{10} = 0.3122 \, \si{rad/s^2} $$

\subsubsection*{3. Cálculo de la posición final}
Se calcula el desplazamiento angular ($\Delta\theta$) para $t = 5 \, \si{s}$:
$$ \Delta\theta = \omega_0 t + \frac{1}{2}\alpha t^2 $$
$$ \Delta\theta = (0.5)(5) + \frac{1}{2}(0.3122)(5^2) $$
$$ \Delta\theta = 2.5 + 3.9025 = 6.4025 \, \si{rad} $$

En grados, esto equivale a:
$$ \Delta\theta = 6.4025 \left(\frac{180^\circ}{\pi}\right) \approx 366.83^\circ $$

La posición angular final ($\theta_f$) es:
$$ \theta_f = \theta_0 + \Delta\theta = 216.87^\circ + 366.83^\circ = 583.7^\circ $$
Restando una vuelta completa ($360^\circ$) se obtiene el ángulo equivalente:
$$ \theta_f = 583.7^\circ - 360^\circ = 223.7^\circ $$

El vector posición final es:
$$ \vec{r}_f = R\cos(223.7^\circ)\vec{i} + R\sin(223.7^\circ)\vec{j} $$
$$ \vec{r}_f = 10(-0.723)\vec{i} + 10(-0.691)\vec{j} $$
$$ \vec{r}_f = -7.23\vec{i} - 6.91\vec{j} \, \si{m} $$

\begin{tcolorbox}[colback=green!5!white, colframe=green!50!black, title=Respuesta Final]
\centering \Large $\vec{r} = -7.22\vec{i} - 6.91\vec{j} \, \si{m}$
\end{tcolorbox}

\newpage

\subsection*{Enunciado}
La partícula de la figura se traslada con MCUVA en sentido antihorario, por una trayectoria de $10 \, \si{m}$ de radio. El módulo de su aceleración inicial fue $4 \, \si{m/s^2}$ y su velocidad inicial $\vec{v}_0 = 3\vec{i} - 4\vec{j} \, \si{m/s}$. Para $t = 5 \, \si{s}$, determine la velocidad de la partícula.

\begin{center}
\begin{tikzpicture}[scale=0.8, >=Latex]
    \draw[->, thick] (-3,0) -- (3,0) node[right] {$x$};
    \draw[->, thick] (0,-3) -- (0,3) node[above] {$y$};
    
    \draw[thick] (0,0) circle (2);
    
    \fill (-1.6,-1.2) circle (0.08);
    \draw[thick, ->] (0,0) -- (-1.6,-1.2);
    
    \draw[thick, ->] (-1.6,-1.2) -- (-0.6,-2.6) node[below] {$\vec{v}_0$};
\end{tikzpicture}
\end{center}

\subsection*{Solución}

\subsubsection*{1. Parámetros previamente calculados}
A partir del análisis del movimiento, se determinaron los siguientes valores constantes y para $t = 5 \, \si{s}$:
$$ \vec{r}_f = -7.23\vec{i} - 6.91\vec{j} \, \si{m} $$
$$ \omega_0 = 0.5 \, \si{rad/s} $$
$$ \alpha = 0.3122 \, \si{rad/s^2} $$

\subsubsection*{2. Cálculo de la velocidad angular final}
Se calcula la velocidad angular ($\omega_f$) a los $5 \, \si{s}$:
$$ \omega_f = \omega_0 + \alpha t $$
$$ \omega_f = 0.5 + (0.3122)(5) $$
$$ \omega_f = 0.5 + 1.561 = 2.061 \, \si{rad/s} $$

Dado que el movimiento es antihorario, el vector velocidad angular final es:
$$ \vec{\omega}_f = 2.061\vec{k} \, \si{rad/s} $$

\subsubsection*{3. Cálculo del vector velocidad final}
Se utiliza la definición de velocidad tangencial en forma vectorial ($\vec{v}_f = \vec{\omega}_f \times \vec{r}_f$):
$$ \vec{v}_f = (2.061\vec{k}) \times (-7.23\vec{i} - 6.91\vec{j}) $$
$$ \vec{v}_f = 2.061(-7.23\vec{j} + 6.91\vec{i}) $$
$$ \vec{v}_f = 14.24\vec{i} - 14.90\vec{j} \, \si{m/s} $$

\begin{tcolorbox}[colback=green!5!white, colframe=green!50!black, title=Respuesta Final]
\centering \Large $\vec{v} = 14.25\vec{i} - 14.89\vec{j} \, \si{m/s}$
\end{tcolorbox}

\newpage

\subsection*{Enunciado}
Una piedra atada a una cuerda gira en el plano horizontal $xz$, con centro en el origen de coordenadas y radio $0.5 \, \si{m}$, a razón de $100 \, \si{rpm}$ en sentido antihorario. En el instante en que la piedra se encuentra en $0.4\vec{i} + 0.3\vec{k} \, \si{m}$, determine la aceleración de la piedra.

\subsection*{Solución}

\subsubsection*{1. Análisis de los datos iniciales}
Se identifican los parámetros dados:
\begin{itemize}
    \item Radio de la trayectoria: $R = 0.5 \, \si{m}$
    \item Rapidez angular: $\omega = 100 \, \si{rpm}$
    \item Vector posición: $\vec{r} = 0.4\vec{i} + 0.3\vec{k} \, \si{m}$
\end{itemize}

Se verifica que el módulo del vector posición coincide con el radio:
$$ |\vec{r}| = \sqrt{0.4^2 + 0.3^2} = \sqrt{0.16 + 0.09} = \sqrt{0.25} = 0.5 \, \si{m} $$

\subsubsection*{2. Conversión de unidades}
Se transforma la rapidez angular de revoluciones por minuto a radianes por segundo:
$$ \omega = 100 \, \frac{\si{rev}}{\si{min}} \times \left( \frac{2\pi \, \si{rad}}{1 \, \si{rev}} \right) \times \left( \frac{1 \, \si{min}}{60 \, \si{s}} \right) $$
$$ \omega = \frac{200\pi}{60} = \frac{10\pi}{3} \, \si{rad/s} $$
Aproximando el valor:
$$ \omega \approx 10.47 \, \si{rad/s} $$

\subsubsection*{3. Cálculo de la aceleración}
En un movimiento circular uniforme (MCU), la única aceleración presente es la aceleración centrípeta ($\vec{a}_c$), la cual siempre apunta hacia el centro de rotación (origen de coordenadas).

Primero, se calcula la magnitud de la aceleración centrípeta:
$$ a_c = \omega^2 R $$
$$ a_c = (10.47)^2 (0.5) $$
$$ a_c \approx 109.62 \times 0.5 = 54.81 \, \si{m/s^2} $$

El vector unitario de la posición ($\vec{u}_r$) se obtiene dividiendo el vector posición por su módulo:
$$ \vec{u}_r = \frac{\vec{r}}{|\vec{r}|} = \frac{0.4\vec{i} + 0.3\vec{k}}{0.5} $$
$$ \vec{u}_r = 0.8\vec{i} + 0.6\vec{k} $$

El vector de aceleración centrípeta se opone al vector posición (apunta hacia el centro):
$$ \vec{a} = \vec{a}_c = -a_c \vec{u}_r $$
$$ \vec{a} = -54.81 (0.8\vec{i} + 0.6\vec{k}) $$
$$ \vec{a} = -43.848\vec{i} - 32.886\vec{k} \, \si{m/s^2} $$

\begin{tcolorbox}[colback=green!5!white, colframe=green!50!black, title=Respuesta Final]
\centering \Large $\vec{a} = -43.85\vec{i} - 32.89\vec{k} \, \si{m/s^2}$
\end{tcolorbox}

\newpage

\subsection*{Enunciado}
Una piedra atada a una cuerda gira en el plano horizontal $xz$, con centro en el origen de coordenadas y radio $0.5 \, \si{m}$, a razón de $100 \, \si{rpm}$ en sentido antihorario. En el instante en que la piedra se encuentra en $0.4\vec{i} + 0.3\vec{k} \, \si{m}$, determine la velocidad angular de la cuerda.

\subsection*{Solución}

\subsubsection*{1. Análisis de la rapidez angular}
Del análisis previo, se determinó que la magnitud de la velocidad angular ($\omega$) en el Sistema Internacional es:
$$ \omega = 100 \, \si{rpm} \times \left( \frac{2\pi \, \si{rad}}{60 \, \si{s}} \right) = \frac{10\pi}{3} \, \si{rad/s} \approx 10.47 \, \si{rad/s} $$

\subsubsection*{2. Determinación de la dirección del vector}
El problema indica que la piedra gira en el plano horizontal $xz$.
La dirección del vector velocidad angular es perpendicular al plano de rotación, de acuerdo con la regla de la mano derecha.
Si el movimiento es en el plano $xz$ y en sentido antihorario, el pulgar apuntará en la dirección positiva del eje restante, es decir, el eje $+y$.

Por lo tanto, el vector unitario que define la dirección de la velocidad angular es $\vec{j}$.

$$ \vec{\omega} = \omega \vec{j} $$
$$ \vec{\omega} = 10.47\vec{j} \, \si{rad/s} $$

\begin{tcolorbox}[colback=green!5!white, colframe=green!50!black, title=Respuesta Final]
\centering \Large $\vec{\omega} = 10.47\vec{j} \, \si{rad/s}$
\end{tcolorbox}

\newpage

\subsection*{Enunciado}
Una ruleta vertical de $0.6 \, \si{m}$ de radio empieza a girar con $3\vec{k} \, \si{rad/s}$. La pluma indicadora del número retarda angularmente su movimiento a razón de $-0.2\vec{k} \, \si{rad/s^2}$. Determine el tiempo que le toma a la ruleta detenerse.

\subsection*{Solución}

\subsubsection*{1. Identificación de datos}
Se extraen los datos cinemáticos iniciales del problema:
$$ \vec{\omega}_0 = 3\vec{k} \, \si{rad/s} \implies \omega_0 = 3 \, \si{rad/s} $$
$$ \vec{\alpha} = -0.2\vec{k} \, \si{rad/s^2} \implies \alpha = -0.2 \, \si{rad/s^2} $$
La condición para que la ruleta se detenga es que su velocidad angular final sea nula:
$$ \omega_f = 0 \, \si{rad/s} $$

\subsubsection*{2. Cálculo del tiempo}
Se utiliza la ecuación de la cinemática circular para la velocidad angular:
$$ \omega_f = \omega_0 + \alpha \Delta t $$

Se sustituyen los valores correspondientes:
$$ 0 = 3 + (-0.2) \Delta t $$
$$ 0.2 \Delta t = 3 $$
$$ \Delta t = \frac{3}{0.2} $$
$$ \Delta t = 15 \, \si{s} $$

\begin{tcolorbox}[colback=green!5!white, colframe=green!50!black, title=Respuesta Final]
\centering \Large $\Delta t = 15 \, \si{s}$
\end{tcolorbox}

\newpage

\subsection*{Enunciado}
Una ruleta vertical de $0.6 \, \si{m}$ de radio empieza a girar con $3\vec{k} \, \si{rad/s}$. La pluma indicadora del número retarda angularmente su movimiento a razón de $-0.2\vec{k} \, \si{rad/s^2}$. Determine el número de revoluciones giradas.

\subsection*{Solución}

\subsubsection*{1. Datos y parámetros previos}
Del análisis cinemático anterior se conocen los siguientes parámetros hasta el instante en que la ruleta se detiene:
$$ \omega_0 = 3 \, \si{rad/s} $$
$$ \alpha = -0.2 \, \si{rad/s^2} $$
$$ \Delta t = 15 \, \si{s} $$

\subsubsection*{2. Cálculo del desplazamiento angular}
Se determina el desplazamiento angular total ($\Delta \theta$) utilizando la ecuación de posición para el movimiento circular uniformemente variado:
$$ \Delta \theta = \omega_0 \Delta t + \frac{1}{2} \alpha \Delta t^2 $$

Se sustituyen los valores:
$$ \Delta \theta = (3)(15) + \frac{1}{2}(-0.2)(15)^2 $$
$$ \Delta \theta = 45 - 0.1(225) $$
$$ \Delta \theta = 45 - 22.5 $$
$$ \Delta \theta = 22.5 \, \si{rad} $$

\subsubsection*{3. Conversión a revoluciones}
Se transforma el desplazamiento angular de radianes a revoluciones sabiendo que una revolución equivale a $2\pi$ radianes:
$$ \text{Nro. de revoluciones} = 22.5 \, \si{rad} \times \left( \frac{1 \, \text{rev}}{2\pi \, \si{rad}} \right) $$
$$ \text{Nro. de revoluciones} = \frac{22.5}{2\pi} \, \text{rev} $$
$$ \text{Nro. de revoluciones} \approx 3.58 \, \text{rev} $$

\begin{tcolorbox}[colback=green!5!white, colframe=green!50!black, title=Respuesta Final]
\centering \Large $3.58 \, \text{rev}$
\end{tcolorbox}

\newpage

\subsection*{Enunciado}
Una partícula se mueve con MCUV en el plano $xy$ en sentido antihorario, por una circunferencia de $2 \, \si{m}$ de radio. En el instante que pasa por $A$, la magnitud de su aceleración es $10 \, \si{m/s^2}$ y su rapidez disminuye a razón de $3 \, \si{m/s}$ cada $0.5 \, \si{s}$. En el instante en que invierte su movimiento, determine la posición de la partícula.

\begin{center}
\begin{tikzpicture}[scale=1, >=Latex]
    \draw[->, thick] (-2.5,0) -- (2.5,0) node[right] {$x$};
    \draw[->, thick] (0,-2.5) -- (0,2.5) node[above] {$y$};
    \draw[thick] (0,0) circle (2);
    \node[below left] at (0,-2) {$A$};
    \draw[->, thick] (-0.3,-2.2) arc (260:280:2.2);
\end{tikzpicture}
\end{center}

\subsection*{Solución}

\subsubsection*{1. Análisis de la aceleración tangencial y normal}
Se calcula la aceleración tangencial ($a_t$) considerando el cambio de rapidez:
$$ a_t = \frac{\Delta v}{\Delta t} = \frac{-3}{0.5} = -6 \, \si{m/s^2} $$

A partir de la magnitud de la aceleración total ($a = 10 \, \si{m/s^2}$), se halla la aceleración centrípeta o normal ($a_n$):
$$ a^2 = a_t^2 + a_n^2 $$
$$ 10^2 = (-6)^2 + a_n^2 $$
$$ 100 = 36 + a_n^2 $$
$$ a_n = \sqrt{64} = 8 \, \si{m/s^2} $$

\subsubsection*{2. Análisis de la cinemática angular}
Se determina la velocidad angular inicial ($\omega_0$) en el punto A:
$$ a_n = \omega_0^2 R $$
$$ 8 = \omega_0^2 (2) $$
$$ \omega_0 = \sqrt{4} = 2 \, \si{rad/s} $$

Se halla la aceleración angular ($\alpha$):
$$ \alpha = \frac{a_t}{R} = \frac{-6}{2} = -3 \, \si{rad/s^2} $$

Cuando la partícula invierte su movimiento, su velocidad angular final es nula ($\omega_f = 0$). Se calcula el desplazamiento angular ($\Delta\theta$):
$$ \omega_f^2 = \omega_0^2 + 2\alpha \Delta\theta $$
$$ 0 = 2^2 + 2(-3)\Delta\theta $$
$$ 0 = 4 - 6\Delta\theta $$
$$ \Delta\theta = \frac{4}{6} = \frac{2}{3} \, \si{rad} \approx 38.2^\circ $$

\subsubsection*{3. Cálculo de la posición}
Considerando que el punto A se encuentra en la parte inferior del eje y, su ángulo inicial es:
$$ \theta_0 = 270^\circ $$

El ángulo de la posición final ($\theta_f$) es:
$$ \theta_f = 270^\circ + 38.2^\circ = 308.2^\circ $$

Se determinan las componentes del vector posición final ($\vec{r}_B$):
$$ \vec{r}_B = R\cos(\theta_f)\vec{i} + R\sin(\theta_f)\vec{j} $$
$$ \vec{r}_B = 2\cos(308.2^\circ)\vec{i} + 2\sin(308.2^\circ)\vec{j} $$
$$ \vec{r}_B = 2(0.6184)\vec{i} + 2(-0.7858)\vec{j} $$
$$ \vec{r}_B = 1.237\vec{i} - 1.572\vec{j} \, \si{m} $$

\begin{tcolorbox}[colback=green!5!white, colframe=green!50!black, title=Respuesta Final]
\centering \Large $\vec{r}_B = 1.24\vec{i} - 1.57\vec{j} \, \si{m}$
\end{tcolorbox}

\newpage

\subsection*{Enunciado}
Una partícula se mueve con MCUV en el plano $xy$ en sentido antihorario, por una circunferencia de $2 \, \si{m}$ de radio. En el instante que pasa por $A$, la magnitud de su aceleración es $10 \, \si{m/s^2}$ y su rapidez disminuye a razón de $3 \, \si{m/s}$ cada $0.5 \, \si{s}$. En el instante en que invierte su movimiento, determinela aceleración de la partícula.

\begin{center}
\begin{tikzpicture}[scale=1, >=Latex]
    \draw[->, thick] (-2.5,0) -- (2.5,0) node[right] {$x$};
    \draw[->, thick] (0,-2.5) -- (0,2.5) node[above] {$y$};
    \draw[thick] (0,0) circle (2);
    \node[below left] at (0,-2) {$A$};
    \draw[->, thick] (-0.3,-2.2) arc (260:280:2.2);
\end{tikzpicture}
\end{center}

\subsection*{Solución}

\subsubsection*{1. Análisis de la aceleración final}
En el instante en que la partícula invierte su movimiento, su velocidad angular es cero ($\omega_f = 0 \, \si{rad/s}$).
Por lo tanto, la aceleración normal o centrípeta es nula:
$$ a_{nf} = \omega_f^2 R = 0 \, \si{m/s^2} $$
La aceleración total de la partícula corresponde únicamente a la aceleración tangencial constante que posee:
$$ \vec{a}_B = \vec{a}_t $$

\subsubsection*{2. Vector aceleración tangencial}
El vector unitario tangencial ($\vec{u}_t$) para un movimiento en sentido antihorario se define como:
$$ \vec{u}_t = -\sin(\theta_f)\vec{i} + \cos(\theta_f)\vec{j} $$

Se evalúa con el ángulo final obtenido ($\theta_f = 308.2^\circ$):
$$ \vec{u}_t = -\sin(308.2^\circ)\vec{i} + \cos(308.2^\circ)\vec{j} $$
$$ \vec{u}_t = -(-0.7858)\vec{i} + 0.6184\vec{j} $$
$$ \vec{u}_t = 0.7858\vec{i} + 0.6184\vec{j} $$

Finalmente, se multiplica el escalar de la aceleración tangencial ($a_t = -6 \, \si{m/s^2}$) por su vector unitario:
$$ \vec{a}_B = a_t \vec{u}_t $$
$$ \vec{a}_B = -6(0.7858\vec{i} + 0.6184\vec{j}) $$
$$ \vec{a}_B = -4.715\vec{i} - 3.710\vec{j} \, \si{m/s^2} $$

\begin{tcolorbox}[colback=green!5!white, colframe=green!50!black, title=Respuesta Final]
\centering \Large $\vec{a}_B = -4.71\vec{i} - 3.71\vec{j} \, \si{m/s^2}$
\end{tcolorbox}

\newpage

\subsection*{Enunciado}
Una partícula se mueve de acuerdo con el gráfico con MCUV, si $R = 10 \, \si{m}$, $v_0 = 20 \, \si{m/s}$, $v_2 = 60 \, \si{m/s}$. Determine la aceleración de la partícula a los $2 \, \si{s}$. 

\begin{center}
\begin{tikzpicture}[scale=1.2, >=Latex]
    \draw[->, thick] (-2.5,0) -- (2.5,0) node[right] {$x$};
    \draw[->, thick] (0,-2.5) -- (0,2.5) node[above] {$y$};
    \draw[thick] (0,0) circle (2);
    
    \fill (0,2) circle (0.06);
    \node[above left] at (0,2) {$t=0$};
    \draw[->, thick] (0,2) -- (1,2) node[above] {$v_0$};
    
    \coordinate (P) at ({2*cos(45)}, {2*sin(45)});
    \fill (P) circle (0.06);
    \node[above right] at (P) {$t=2 \, \si{s}$};
    
    \draw[->, thick] (0,0) -- (P);
    \draw (0.5,0) arc (0:45:0.5);
    \node at (0.8, 0.25) {\small $45^\circ$};
    
    \draw[->, thick] (P) -- ++({1.2*sin(45)}, {-1.2*cos(45)}) node[right] {$v$};
\end{tikzpicture}
\end{center}

\subsection*{Solución}

\subsubsection*{1. Análisis de magnitudes cinemáticas}
Se determinan las velocidades angulares inicial ($t=0$) y final ($t=2 \, \si{s}$):
$$ \omega_0 = \frac{v_0}{R} = \frac{20}{10} = 2 \, \si{rad/s} $$
$$ \omega_f = \frac{v_2}{R} = \frac{60}{10} = 6 \, \si{rad/s} $$

Se calcula la aceleración angular ($\alpha$):
$$ \alpha = \frac{\omega_f - \omega_0}{\Delta t} = \frac{6 - 2}{2} = 2 \, \si{rad/s^2} $$

Con esto, se hallan las magnitudes de la aceleración tangencial ($a_t$) y normal ($a_n$) a los $2 \, \si{s}$:
$$ a_t = \alpha R = 2(10) = 20 \, \si{m/s^2} $$
$$ a_n = \omega_f^2 R = 6^2(10) = 360 \, \si{m/s^2} $$

\subsubsection*{2. Análisis vectorial en t = 2 s}
Según el gráfico, a los $2 \, \si{s}$ la partícula se encuentra en la posición angular de $45^\circ$ y se mueve en sentido horario.

El vector aceleración normal ($\vec{a}_n$) apunta hacia el centro de la trayectoria (origen):
$$ \vec{u}_n = -\cos(45^\circ)\vec{i} - \sin(45^\circ)\vec{j} = -\frac{\sqrt{2}}{2}\vec{i} - \frac{\sqrt{2}}{2}\vec{j} $$
$$ \vec{a}_n = 360 \left(-\frac{\sqrt{2}}{2}\vec{i} - \frac{\sqrt{2}}{2}\vec{j}\right) = -180\sqrt{2}\vec{i} - 180\sqrt{2}\vec{j} \, \si{m/s^2} $$

El vector aceleración tangencial ($\vec{a}_t$) es paralelo a la velocidad. Para un movimiento horario en el primer cuadrante, el vector unitario tangencial es:
$$ \vec{u}_t = \sin(45^\circ)\vec{i} - \cos(45^\circ)\vec{j} = \frac{\sqrt{2}}{2}\vec{i} - \frac{\sqrt{2}}{2}\vec{j} $$
$$ \vec{a}_t = 20 \left(\frac{\sqrt{2}}{2}\vec{i} - \frac{\sqrt{2}}{2}\vec{j}\right) = 10\sqrt{2}\vec{i} - 10\sqrt{2}\vec{j} \, \si{m/s^2} $$

\subsubsection*{3. Aceleración total}
Se suman las componentes vectoriales:
$$ \vec{a} = \vec{a}_n + \vec{a}_t $$
$$ \vec{a} = (-180\sqrt{2} + 10\sqrt{2})\vec{i} + (-180\sqrt{2} - 10\sqrt{2})\vec{j} $$
$$ \vec{a} = -170\sqrt{2}\vec{i} - 190\sqrt{2}\vec{j} \, \si{m/s^2} $$

Aproximando el valor de $\sqrt{2} \approx 1.4142$:
$$ \vec{a} = -170(1.4142)\vec{i} - 190(1.4142)\vec{j} $$
$$ \vec{a} = -240.414\vec{i} - 268.698\vec{j} \, \si{m/s^2} $$

\begin{tcolorbox}[colback=green!5!white, colframe=green!50!black, title=Respuesta Final]
\centering \Large $\vec{a} = -240.42\vec{i} - 268.7\vec{j} \, \si{m/s^2}$
\end{tcolorbox}

\newpage

\subsection*{Enunciado}
Una partícula se mueve a lo largo de la trayectoria circular de $2 \, \si{m}$ de radio, su vector posición tiene una aceleración angular constante $\vec{\alpha} = 2\vec{k} \, \si{rad/s^2}$. Al instante $t = 0 \, \si{s}$ su vector posición tiene una velocidad angular constante $\vec{\omega} = -10\vec{k} \, \si{rad/s}$. Si a los $2 \, \si{s}$ pasa por $A$, determine la aceleración centrípeta de la partícula en ese instante. 

\begin{center}
\begin{tikzpicture}[scale=1.2, >=Latex]
    \draw[->, thick] (-2.5,0) -- (2.5,0) node[right] {$x$};
    \draw[->, thick] (0,-2.5) -- (0,2.5) node[above] {$y$};
    \draw[thick] (0,0) circle (2);
    
    \coordinate (A) at ({2*cos(-30)}, {2*sin(-30)});
    \fill (A) circle (0.06);
    \node[below right] at (A) {$A$};
    
    \draw[->, thick] (0,0) -- (A);
    
    \draw (0.6,0) arc (0:-30:0.6);
    \node at (0.9, -0.25) {\small $30^\circ$};
\end{tikzpicture}
\end{center}

\subsection*{Solución}

\subsubsection*{1. Cálculo de la velocidad angular en t = 2 s}
A partir del enunciado, se identifican la velocidad angular inicial ($\vec{\omega}_0$) en $t = 0 \, \si{s}$ y la aceleración angular constante ($\vec{\alpha}$):
$$ \vec{\omega}_0 = -10\vec{k} \, \si{rad/s} $$
$$ \vec{\alpha} = 2\vec{k} \, \si{rad/s^2} $$

Se utiliza la ecuación de la cinemática circular para determinar la velocidad angular final ($\vec{\omega}_f$) a los $2 \, \si{s}$:
$$ \vec{\omega}_f = \vec{\omega}_0 + \vec{\alpha} t $$
$$ \vec{\omega}_f = -10\vec{k} + (2\vec{k})(2) $$
$$ \vec{\omega}_f = -10\vec{k} + 4\vec{k} $$
$$ \vec{\omega}_f = -6\vec{k} \, \si{rad/s} $$

\subsubsection*{2. Determinación del vector posición}
En el instante $t = 2 \, \si{s}$, la partícula se ubica en el punto $A$. Según el gráfico, el vector posición forma un ángulo de $30^\circ$ por debajo del eje positivo de las abscisas (cuarto cuadrante).

El módulo del vector posición es igual al radio de la trayectoria circular ($R = 2 \, \si{m}$). Se determinan sus componentes rectangulares:
$$ \vec{r} = R\cos(-30^\circ)\vec{i} + R\sin(-30^\circ)\vec{j} $$
$$ \vec{r} = 2\left(\frac{\sqrt{3}}{2}\right)\vec{i} + 2\left(-\frac{1}{2}\right)\vec{j} $$
$$ \vec{r} = \sqrt{3}\vec{i} - \vec{j} \, \si{m} $$

\subsubsection*{3. Cálculo de la aceleración centrípeta}
La aceleración centrípeta ($\vec{a}_c$) se obtiene vectorialmente mediante la relación $\vec{a}_c = -\omega_f^2 \vec{r}$, donde $\omega_f$ es la magnitud de la velocidad angular en ese instante ($\omega_f = |-6\vec{k}| = 6 \, \si{rad/s}$).

Sustituyendo los valores calculados:
$$ \vec{a}_c = -(6)^2 (\sqrt{3}\vec{i} - \vec{j}) $$
$$ \vec{a}_c = -36 (\sqrt{3}\vec{i} - \vec{j}) $$
$$ \vec{a}_c = -36\sqrt{3}\vec{i} + 36\vec{j} \, \si{m/s^2} $$

Aproximando el valor de la raíz cuadrada ($\sqrt{3} \approx 1.73205$):
$$ \vec{a}_c = -36(1.73205)\vec{i} + 36\vec{j} \, \si{m/s^2} $$
$$ \vec{a}_c = -62.3538\vec{i} + 36\vec{j} \, \si{m/s^2} $$

\begin{tcolorbox}[colback=green!5!white, colframe=green!50!black, title=Respuesta Final]
\centering \Large $\vec{a}_c = -62.35\vec{i} + 36\vec{j} \, \si{m/s^2}$
\end{tcolorbox}

\newpage

\subsection*{Enunciado}
Una partícula se mueve antihorariamente a lo largo de una circunferencia de $1.2 \, \si{m}$ de radio en el plano $xy$. En el instante $0 \, \si{s}$ su rapidez es $2 \, \si{m/s}$, la posición angular inicial de su radio vector es $\pi/6 \, \si{rad}$ con una aceleración angular constante de $2\vec{k} \, \si{rad/s^2}$. Determine la aceleración de la partícula a los $2 \, \si{s}$.
\subsection*{Solución}

\subsubsection*{1. Cálculo de parámetros angulares}
Se halla la velocidad angular inicial ($\omega_0$):
$$ \omega_0 = \frac{v_0}{R} = \frac{2}{1.2} = \frac{5}{3} \, \si{rad/s} $$

Se determina la velocidad angular a los $2 \, \si{s}$ ($\omega_f$):
$$ \omega_f = \omega_0 + \alpha t $$
$$ \omega_f = \frac{5}{3} + (2)(2) = \frac{17}{3} \, \si{rad/s} \approx 5.667 \, \si{rad/s} $$

Se calcula la posición angular final ($\theta_f$):
$$ \theta_f = \theta_0 + \omega_0 t + \frac{1}{2}\alpha t^2 $$
$$ \theta_f = \frac{\pi}{6} + \left(\frac{5}{3}\right)(2) + \frac{1}{2}(2)(2)^2 $$
$$ \theta_f = \frac{\pi}{6} + \frac{10}{3} + 4 \approx 0.5236 + 3.3333 + 4 = 7.8569 \, \si{rad} $$

Convirtiendo a grados, esto equivale a $450.17^\circ$. Restando una vuelta completa ($360^\circ$), la posición en la primera revolución es aproximadamente $90.17^\circ$.

\subsubsection*{2. Vector posición}
Se hallan las componentes rectangulares del vector posición ($\vec{r}$) a los $2 \, \si{s}$:
$$ \vec{r} = R\cos(\theta_f)\vec{i} + R\sin(\theta_f)\vec{j} $$
$$ \vec{r} = 1.2\cos(7.8569 \, \si{rad})\vec{i} + 1.2\sin(7.8569 \, \si{rad})\vec{j} $$
$$ \vec{r} = 1.2(-0.0035)\vec{i} + 1.2(0.9999)\vec{j} $$
$$ \vec{r} = -0.0042\vec{i} + 1.2\vec{j} \, \si{m} $$

\subsubsection*{3. Cálculo de la aceleración total}
Se calcula la aceleración normal o centrípeta ($\vec{a}_n$):
$$ \vec{a}_n = -\omega_f^2 \vec{r} $$
$$ \vec{a}_n = -\left(\frac{17}{3}\right)^2 (-0.0042\vec{i} + 1.2\vec{j}) $$
$$ \vec{a}_n = -32.111 (-0.0042\vec{i} + 1.2\vec{j}) $$
$$ \vec{a}_n = 0.135\vec{i} - 38.533\vec{j} \, \si{m/s^2} $$

Se halla la aceleración tangencial ($\vec{a}_t$) usando el producto cruz entre el vector aceleración angular y la posición:
$$ \vec{a}_t = \vec{\alpha} \times \vec{r} $$
$$ \vec{a}_t = (2\vec{k}) \times (-0.0042\vec{i} + 1.2\vec{j}) $$
$$ \vec{a}_t = -2.4\vec{i} - 0.008\vec{j} \, \si{m/s^2} $$

Finalmente, se suman las componentes para obtener la aceleración total ($\vec{a}$):
$$ \vec{a} = \vec{a}_n + \vec{a}_t $$
$$ \vec{a} = (0.135 - 2.4)\vec{i} + (-38.533 - 0.008)\vec{j} $$
$$ \vec{a} = -2.265\vec{i} - 38.541\vec{j} \, \si{m/s^2} $$

*(Nota: La respuesta calculada numéricamente de forma exacta difiere muy levemente de la proporcionada en el libro base debido a que dicho texto probablemente aproximó la posición final a $90^\circ$ exactos al plantearlo)*.

\begin{tcolorbox}[colback=green!5!white, colframe=green!50!black, title=Respuesta Final]
\centering \Large $\vec{a} = -2.27\vec{i} - 38.54\vec{j} \, \si{m/s^2}$
\end{tcolorbox}

\newpage

\subsection*{Enunciado}
Una partícula se mueve con MCUV en el plano $xy$. Si a los $0 \, \si{s}$, se encuentra en la posición $(-3\vec{i} - 4\vec{j}) \, \si{m}$, con aceleración $(8\vec{i} + 3\vec{j}) \, \si{m/s^2}$, determine la aceleración de la partícula cuando $t = 2 \, \si{s}$. 

\begin{center}
\begin{tikzpicture}[scale=0.8, >=Latex]
    \draw[->, thick] (-3.5,0) -- (3.5,0) node[right] {$x$};
    \draw[->, thick] (0,-3.5) -- (0,3.5) node[above] {$y$};
    \draw[thick] (0,0) circle (2.5);
    
    \coordinate (P) at (-1.5, -2);
    \draw[->, thick] (0,0) -- (P) node[midway, below right] {$\vec{r}_0$};
    
    \draw[->, thick, blue] (P) -- ++(1.5, 2) node[right] {$\vec{a}_n$};
    \draw[->, thick, blue] (P) -- ++(-1.6, 1.2) node[left] {$\vec{a}_t$};
    \draw[->, thick, red] (P) -- ++(-0.1, 3.2) node[above right] {$\vec{a}$};
\end{tikzpicture}
\end{center}

\subsection*{Solución}

\textit{Nota: La respuesta proporcionada por el libro base es incorrecta. A continuación se presenta el desarrollo analítico que demuestra la verdadera aceleración de la partícula.}

\subsubsection*{1. Análisis de la posición y dirección inicial}
El vector posición inicial es $\vec{r}_0 = -3\vec{i} - 4\vec{j} \, \si{m}$.
Se calcula el radio de la trayectoria:
$$ R = |\vec{r}_0| = \sqrt{(-3)^2 + (-4)^2} = \sqrt{9 + 16} = 5 \, \si{m} $$

El vector unitario radial en la posición inicial es:
$$ \vec{u}_{r0} = \frac{\vec{r}_0}{R} = -0.6\vec{i} - 0.8\vec{j} $$

El vector unitario normal (dirigido hacia el centro) es opuesto al radial:
$$ \vec{u}_{n0} = -\vec{u}_{r0} = 0.6\vec{i} + 0.8\vec{j} $$

\subsubsection*{2. Componentes de la aceleración inicial}
La aceleración total inicial es $\vec{a}_0 = 8\vec{i} + 3\vec{j} \, \si{m/s^2}$.
Se halla la magnitud de la aceleración normal ($a_{n0}$) proyectando $\vec{a}_0$ sobre $\vec{u}_{n0}$ mediante el producto punto:
$$ a_{n0} = \vec{a}_0 \cdot \vec{u}_{n0} = (8)(0.6) + (3)(0.8) = 4.8 + 2.4 = 7.2 \, \si{m/s^2} $$

Con esto, se obtiene el vector aceleración normal inicial:
$$ \vec{a}_{n0} = a_{n0}\vec{u}_{n0} = 7.2(0.6\vec{i} + 0.8\vec{j}) = 4.32\vec{i} + 5.76\vec{j} \, \si{m/s^2} $$

El vector aceleración tangencial inicial ($\vec{a}_{t0}$) se obtiene por diferencia ($\vec{a}_0 = \vec{a}_{n0} + \vec{a}_{t0}$):
$$ \vec{a}_{t0} = \vec{a}_0 - \vec{a}_{n0} = (8\vec{i} + 3\vec{j}) - (4.32\vec{i} + 5.76\vec{j}) = 3.68\vec{i} - 2.76\vec{j} \, \si{m/s^2} $$

La magnitud de la aceleración tangencial constante es:
$$ a_t = \sqrt{(3.68)^2 + (-2.76)^2} = \sqrt{13.5424 + 7.6176} = \sqrt{21.16} = 4.6 \, \si{m/s^2} $$

\subsubsection*{3. Parámetros angulares}
Se determina la velocidad angular inicial ($\omega_0$):
$$ a_{n0} = \omega_0^2 R \implies \omega_0 = \sqrt{\frac{a_{n0}}{R}} = \sqrt{\frac{7.2}{5}} = \sqrt{1.44} = 1.2 \, \si{rad/s} $$

Se calcula la aceleración angular constante ($\alpha$):
$$ \alpha = \frac{a_t}{R} = \frac{4.6}{5} = 0.92 \, \si{rad/s^2} $$

El ángulo de la posición inicial ($\theta_0$) en el tercer cuadrante es:
$$ \theta_0 = 180^\circ + \arctan\left(\frac{-4}{-3}\right) = 180^\circ + 53.13^\circ = 233.13^\circ \approx 4.069 \, \si{rad} $$

\subsubsection*{4. Cinemática a los t = 2 s}
Se halla la velocidad angular final ($\omega_f$):
$$ \omega_f = \omega_0 + \alpha t = 1.2 + (0.92)(2) = 1.2 + 1.84 = 3.04 \, \si{rad/s} $$

Se calcula la posición angular final ($\theta_f$):
$$ \theta_f = \theta_0 + \omega_0 t + \frac{1}{2}\alpha t^2 $$
$$ \theta_f = 4.069 + 1.2(2) + \frac{1}{2}(0.92)(2^2) $$
$$ \theta_f = 4.069 + 2.4 + 1.84 = 8.309 \, \si{rad} $$
Restando una revolución ($2\pi \approx 6.283 \, \si{rad}$), la posición equivale a $2.026 \, \si{rad}$ o $116.08^\circ$ (segundo cuadrante).

\subsubsection*{5. Aceleración final a los t = 2 s}
Se hallan los vectores unitarios radial ($\vec{u}_{rf}$) y tangencial ($\vec{u}_{tf}$) para la nueva posición:
$$ \vec{u}_{rf} = \cos(116.08^\circ)\vec{i} + \sin(116.08^\circ)\vec{j} = -0.440\vec{i} + 0.898\vec{j} $$
$$ \vec{u}_{tf} = -\sin(116.08^\circ)\vec{i} + \cos(116.08^\circ)\vec{j} = -0.898\vec{i} - 0.440\vec{j} $$

Se calcula el vector aceleración normal final ($\vec{a}_{nf}$):
$$ a_{nf} = \omega_f^2 R = (3.04)^2 (5) = 46.208 \, \si{m/s^2} $$
$$ \vec{a}_{nf} = -a_{nf}\vec{u}_{rf} = -46.208(-0.440\vec{i} + 0.898\vec{j}) = 20.33\vec{i} - 41.49\vec{j} \, \si{m/s^2} $$

El vector aceleración tangencial final ($\vec{a}_{tf}$) es:
$$ \vec{a}_{tf} = a_t \vec{u}_{tf} = 4.6(-0.898\vec{i} - 0.440\vec{j}) = -4.13\vec{i} - 2.02\vec{j} \, \si{m/s^2} $$

Finalmente, la aceleración total ($\vec{a}_f$) es la suma de sus componentes:
$$ \vec{a}_f = \vec{a}_{nf} + \vec{a}_{tf} $$
$$ \vec{a}_f = (20.33 - 4.13)\vec{i} + (-41.49 - 2.02)\vec{j} $$
$$ \vec{a}_f = 16.20\vec{i} - 43.51\vec{j} \, \si{m/s^2} $$

\begin{tcolorbox}[colback=green!5!white, colframe=green!50!black, title=Respuesta Final]
\centering \Large $\vec{a} = 16.20\vec{i} - 43.51\vec{j} \, \si{m/s^2}$
\end{tcolorbox}

\newpage

\subsection*{Enunciado}
Dos autos $A$ y $B$ viajan sobre una redondel horizontal de $30 \, \si{m}$ de radio en el plano $xy$, considerando a los autos como partículas y que se trasladan en sentidos contrarios aceleradamente, con $v_{A_0} = 12 \, \si{m/s}$ y $v_{B_0} = 17 \, \si{m/s}$, $a_{T_A} = a_{T_B} = 0.1 \, \si{m/s^2}$, cte. Luego de $5 \, \si{s}$ de la posición mostrada, determine la posición relativa de $A$ respecto a $B$. 

\begin{center}
\begin{tikzpicture}[scale=1, >=Latex]
    \draw[->, thick] (-2.5,0) -- (2.5,0) node[right] {$x$};
    \draw[->, thick] (0,-2.5) -- (0,2.5) node[above] {$y$};
    \draw[thick] (0,0) circle (2);
    
    \fill (-2,0) circle (0.08);
    \draw[->, thick] (-2,0) -- (-2,1) node[left] {$v_{A_0}$};
    
    \fill (0,-2) circle (0.08);
    \draw[->, thick] (0,-2) -- (1,-2) node[below right] {$v_{B_0}$};
    
    \draw[->, thick] (0,0) -- (-1.414, 1.414) node[midway, below left] {R};
\end{tikzpicture}
\end{center}

\subsection*{Solución}

\subsubsection*{1. Análisis cinemático del auto A}
El auto A se encuentra en la posición inicial $(-30\vec{i}) \, \si{m}$, lo que corresponde a un ángulo inicial $\theta_{A0} = \pi \, \si{rad}$.
Se observa que la velocidad apunta en la dirección $+y$, lo que indica un movimiento en sentido horario. Por tanto, su velocidad angular inicial es negativa:
$$ \omega_{A0} = -\frac{v_{A0}}{R} = -\frac{12}{30} = -0.4 \, \si{rad/s} $$

Como el auto se traslada "aceleradamente", la aceleración angular debe tener el mismo sentido (horario) que la velocidad angular para aumentar su rapidez:
$$ \alpha_A = -\frac{a_{TA}}{R} = -\frac{0.1}{30} = -\frac{1}{300} \, \si{rad/s^2} $$

Se calcula la posición angular de A a los $5 \, \si{s}$:
$$ \theta_A = \theta_{A0} + \omega_{A0} t + \frac{1}{2}\alpha_A t^2 $$
$$ \theta_A = \pi + (-0.4)(5) + \frac{1}{2}\left(-\frac{1}{300}\right)(5^2) $$
$$ \theta_A = \pi - 2 - \frac{25}{600} = \pi - 2 - \frac{1}{24} $$
$$ \theta_A \approx 3.14159 - 2 - 0.04167 = 1.09992 \, \si{rad} $$

Se convierte a grados para el cálculo del vector posición ($1.09992 \times \frac{180^\circ}{\pi} \approx 63.02^\circ$):
$$ \vec{r}_A = 30\cos(63.02^\circ)\vec{i} + 30\sin(63.02^\circ)\vec{j} $$
$$ \vec{r}_A = 30(0.4537)\vec{i} + 30(0.8912)\vec{j} $$
$$ \vec{r}_A = 13.611\vec{i} + 26.736\vec{j} \, \si{m} $$

\subsubsection*{2. Análisis cinemático del auto B}
El auto B está en la posición $( -30\vec{j} ) \, \si{m}$, equivalente a un ángulo inicial $\theta_{B0} = \frac{3\pi}{2} \, \si{rad}$.
Su velocidad apunta en la dirección $+x$, indicando un movimiento en sentido antihorario (positivo):
$$ \omega_{B0} = \frac{v_{B0}}{R} = \frac{17}{30} \approx 0.5667 \, \si{rad/s} $$

De igual manera, al ser un movimiento acelerado, la aceleración angular es positiva:
$$ \alpha_B = \frac{a_{TB}}{R} = \frac{0.1}{30} = \frac{1}{300} \, \si{rad/s^2} $$

Se calcula la posición angular de B a los $5 \, \si{s}$:
$$ \theta_B = \theta_{B0} + \omega_{B0} t + \frac{1}{2}\alpha_B t^2 $$
$$ \theta_B = \frac{3\pi}{2} + \left(\frac{17}{30}\right)(5) + \frac{1}{2}\left(\frac{1}{300}\right)(25) $$
$$ \theta_B = 1.5\pi + \frac{85}{30} + \frac{1}{24} \approx 4.71239 + 2.83333 + 0.04167 = 7.58739 \, \si{rad} $$

Restando una vuelta completa ($2\pi \approx 6.28319 \, \si{rad}$) para obtener el ángulo dentro de la primera revolución:
$$ \theta_B = 7.58739 - 6.28319 = 1.3042 \, \si{rad} $$

En grados, equivale a $1.3042 \times \frac{180^\circ}{\pi} \approx 74.72^\circ$:
$$ \vec{r}_B = 30\cos(74.72^\circ)\vec{i} + 30\sin(74.72^\circ)\vec{j} $$
$$ \vec{r}_B = 30(0.2635)\vec{i} + 30(0.9647)\vec{j} $$
$$ \vec{r}_B = 7.905\vec{i} + 28.941\vec{j} \, \si{m} $$

\subsubsection*{3. Cálculo de la posición relativa}
La posición relativa de A respecto a B ($\vec{r}_{A/B}$) se determina restando el vector posición de B al vector posición de A:
$$ \vec{r}_{A/B} = \vec{r}_A - \vec{r}_B $$
$$ \vec{r}_{A/B} = (13.611\vec{i} + 26.736\vec{j}) - (7.905\vec{i} + 28.941\vec{j}) $$
$$ \vec{r}_{A/B} = (13.611 - 7.905)\vec{i} + (26.736 - 28.941)\vec{j} $$
$$ \vec{r}_{A/B} = 5.706\vec{i} - 2.205\vec{j} \, \si{m} $$

\begin{tcolorbox}[colback=green!5!white, colframe=green!50!black, title=Respuesta Final]
\centering \Large $\vec{r}_{A/B} = 5.71\vec{i} - 2.21\vec{j} \, \si{m}$
\end{tcolorbox}

\newpage

\subsection*{Enunciado}
Dos autos $A$ y $B$ viajan sobre una redondel horizontal de $30 \, \si{m}$ de radio en el plano $xy$, considerando a los autos como partículas y que se trasladan en sentidos contrarios aceleradamente, con $v_{A_0} = 12 \, \si{m/s}$ y $v_{B_0} = 17 \, \si{m/s}$, $a_{T_A} = a_{T_B} = 0.1 \, \si{m/s^2}$, cte. Luego de $5 \, \si{s}$ de la posición mostrada, determinela aceleración angular del vector posición de $A$.

\begin{center}
\begin{tikzpicture}[scale=0.8, >=Latex]
    \draw[->, thick] (-2.5,0) -- (2.5,0) node[right] {$x$};
    \draw[->, thick] (0,-2.5) -- (0,2.5) node[above] {$y$};
    \draw[thick] (0,0) circle (2);
    
    \fill (-2,0) circle (0.08);
    \draw[->, thick] (-2,0) -- (-2,1) node[left] {$v_{A_0}$};
    
    \fill (0,-2) circle (0.08);
    \draw[->, thick] (0,-2) -- (1,-2) node[below right] {$v_{B_0}$};
    
    \draw[->, thick] (0,0) -- (-1.414, 1.414) node[midway, below left] {R};
\end{tikzpicture}
\end{center}

\subsection*{Solución}

\subsubsection*{1. Relación entre aceleración tangencial y angular}
La magnitud de la aceleración angular ($\alpha$) se relaciona con la aceleración tangencial ($a_T$) mediante el radio de curvatura ($R$):
$$ \alpha_A = \frac{a_{TA}}{R} $$

Sustituyendo los valores dados en el enunciado:
$$ \alpha_A = \frac{0.1}{30} = \frac{1}{300} \approx 0.00333 \, \si{rad/s^2} $$
En notación científica, esto es $3.33 \times 10^{-3} \, \si{rad/s^2}$.

\subsubsection*{2. Determinación del sentido vectorial}
Se analiza el movimiento del auto A desde su posición inicial en el eje $-x$.
Dado que su velocidad inicial apunta hacia la dirección $+y$, la partícula se mueve en sentido de las manecillas del reloj (horario).

El problema indica que los autos se trasladan "aceleradamente", lo que significa que la rapidez está aumentando. Por lo tanto, el vector aceleración tangencial tiene la misma dirección y sentido que el vector velocidad. En consecuencia, la aceleración angular acompaña el sentido del movimiento horario.

Vectorialmente, el sentido horario corresponde a la dirección negativa del eje $z$ (vector unitario $-\vec{k}$).
$$ \vec{\alpha}_A = -3.33 \times 10^{-3} \vec{k} \, \si{rad/s^2} $$

\begin{tcolorbox}[colback=green!5!white, colframe=green!50!black, title=Respuesta Final]
\centering \Large $\vec{\alpha}_A = -3.33 \times 10^{-3}\vec{k} \, \si{rad/s^2}$
\end{tcolorbox}

\newpage

\subsection*{Enunciado}
Una partícula con MCUVA se mueve en el plano $xy$ en sentido antihorario durante $5 \, \si{s}$ por una circunferencia de $10 \, \si{m}$ de radio. Al instante $t = 0 \, \si{s}$, la magnitud de su aceleración es $a_0 = 4 \, \si{m/s^2}$ y su velocidad $\vec{v}_0 = 3\vec{i} - 4\vec{j} \, \si{m/s}$. Determine la velocidad media de la partícula en el intervalo de $0$ a $5 \, \si{s}$. 

\subsection*{Solución}

\textit{Nota: La respuesta proporcionada por el libro base es incorrecta. A continuación se presenta el desarrollo analítico exacto.}

\subsubsection*{1. Análisis de condiciones iniciales (t = 0 s)}
Se calcula la magnitud de la velocidad inicial:
$$ v_0 = |\vec{v}_0| = \sqrt{3^2 + (-4)^2} = 5 \, \si{m/s} $$

La velocidad angular inicial ($\omega_0$) es:
$$ \omega_0 = \frac{v_0}{R} = \frac{5}{10} = 0.5 \, \si{rad/s} $$

Para hallar el vector posición inicial ($\vec{r}_0$), se usa la relación del producto cruz $\vec{v}_0 = \vec{\omega}_0 \times \vec{r}_0$. Sabiendo que gira en sentido antihorario ($\vec{\omega}_0 = 0.5\vec{k}$):
$$ 3\vec{i} - 4\vec{j} = (0.5\vec{k}) \times (x\vec{i} + y\vec{j}) $$
$$ 3\vec{i} - 4\vec{j} = -0.5y\vec{i} + 0.5x\vec{j} $$
Igualando componentes:
$$ -0.5y = 3 \implies y = -6 $$
$$ 0.5x = -4 \implies x = -8 $$
$$ \vec{r}_0 = -8\vec{i} - 6\vec{j} \, \si{m} $$

El ángulo de la posición inicial ($\theta_0$) en el tercer cuadrante es:
$$ \theta_0 = 180^\circ + \arctan\left(\frac{-6}{-8}\right) = 180^\circ + 36.87^\circ = 216.87^\circ \approx 3.785 \, \si{rad} $$

\subsubsection*{2. Cálculo de la aceleración angular}
Se determina la aceleración centrípeta inicial ($a_{n0}$):
$$ a_{n0} = \frac{v_0^2}{R} = \frac{5^2}{10} = 2.5 \, \si{m/s^2} $$

A partir de la aceleración total inicial ($a_0 = 4 \, \si{m/s^2}$), se calcula la aceleración tangencial ($a_t$):
$$ a_0^2 = a_{n0}^2 + a_t^2 $$
$$ 4^2 = 2.5^2 + a_t^2 $$
$$ 16 = 6.25 + a_t^2 $$
$$ a_t = \sqrt{9.75} \approx 3.122 \, \si{m/s^2} $$

La aceleración angular constante ($\alpha$) es:
$$ \alpha = \frac{a_t}{R} = \frac{\sqrt{9.75}}{10} \approx 0.3122 \, \si{rad/s^2} $$

\subsubsection*{3. Posición final a los 5 s}
Se calcula el desplazamiento angular ($\Delta\theta$) para $t = 5 \, \si{s}$:
$$ \Delta\theta = \omega_0 t + \frac{1}{2}\alpha t^2 $$
$$ \Delta\theta = (0.5)(5) + \frac{1}{2}(0.3122)(5^2) $$
$$ \Delta\theta = 2.5 + 3.9025 = 6.4025 \, \si{rad} $$

En grados, equivale a:
$$ \Delta\theta = 6.4025 \left(\frac{180^\circ}{\pi}\right) \approx 366.84^\circ $$

La posición angular final ($\theta_f$) es:
$$ \theta_f = \theta_0 + \Delta\theta = 216.87^\circ + 366.84^\circ = 583.71^\circ $$
Restando una vuelta completa ($360^\circ$):
$$ \theta_f = 223.71^\circ $$

El vector posición final ($\vec{r}_f$) es:
$$ \vec{r}_f = 10\cos(223.71^\circ)\vec{i} + 10\sin(223.71^\circ)\vec{j} $$
$$ \vec{r}_f = 10(-0.7228)\vec{i} + 10(-0.6910)\vec{j} $$
$$ \vec{r}_f = -7.228\vec{i} - 6.910\vec{j} \, \si{m} $$

\subsubsection*{4. Cálculo de la velocidad media}
El desplazamiento de la partícula ($\Delta\vec{r}$) es:
$$ \Delta\vec{r} = \vec{r}_f - \vec{r}_0 $$
$$ \Delta\vec{r} = (-7.228\vec{i} - 6.910\vec{j}) - (-8\vec{i} - 6\vec{j}) $$
$$ \Delta\vec{r} = 0.772\vec{i} - 0.910\vec{j} \, \si{m} $$

La velocidad media ($\vec{v}_m$) en el intervalo de $5 \, \si{s}$ se obtiene por definición:
$$ \vec{v}_m = \frac{\Delta\vec{r}}{\Delta t} $$
$$ \vec{v}_m = \frac{0.772\vec{i} - 0.910\vec{j}}{5} $$
$$ \vec{v}_m = 0.1544\vec{i} - 0.1820\vec{j} \, \si{m/s} $$

\begin{tcolorbox}[colback=green!5!white, colframe=green!50!black, title=Respuesta Final]
\centering \Large $\vec{v}_m = 0.16\vec{i} - 0.18\vec{j} \, \si{m/s}$
\end{tcolorbox}

\newpage

\subsection*{Enunciado}
Una partícula se desplaza con MCUVR en el plano $xy$, al instante $t = 0 \, \si{s}$ pasa por el punto $A(4; 3) \, \si{m}$, con rapidez de $6 \, \si{m/s}$, en sentido antihorario y una aceleración tangencial de $3.1 \, \si{m/s^2}$. Luego de $1.5 \, \si{s}$, pasa por primera vez por el punto $B(-5; 0) \, \si{m}$. Determine el número de revoluciones de su radio vector de $0$ a $10 \, \si{s}$. 

\subsection*{Solución}

\textit{Nota: La respuesta proporcionada por el libro base es incorrecta. A continuación se presenta el desarrollo analítico que demuestra el número real de revoluciones.}

\subsubsection*{1. Análisis de las condiciones iniciales}
A partir del punto inicial $A(4; 3) \, \si{m}$, se calcula el radio de la trayectoria circular:
$$ R = \sqrt{4^2 + 3^2} = \sqrt{16 + 9} = \sqrt{25} = 5 \, \si{m} $$

Se halla la magnitud de la velocidad angular inicial ($\omega_0$):
$$ \omega_0 = \frac{v_0}{R} = \frac{6}{5} = 1.2 \, \si{rad/s} $$
Como el movimiento inicial es en sentido antihorario, el vector es $\vec{\omega}_0 = 1.2\vec{k} \, \si{rad/s}$.

\subsubsection*{2. Aceleración angular}
Se calcula la magnitud de la aceleración angular ($\alpha$) a partir de la aceleración tangencial:
$$ \alpha = \frac{a_t}{R} = \frac{3.1}{5} = 0.62 \, \si{rad/s^2} $$

El problema indica que es un Movimiento Circular Uniformemente Variado Retardado (MCUVR). Al ser retardado, la aceleración angular debe oponerse a la velocidad angular. Por lo tanto, actúa en sentido horario:
$$ \vec{\alpha} = -0.62\vec{k} \, \si{rad/s^2} $$

\subsubsection*{3. Cálculo del desplazamiento angular y revoluciones}
Se evalúa el desplazamiento angular ($\Delta\theta$) para el intervalo $t = 10 \, \si{s}$:
$$ \Delta\theta = \omega_0 t + \frac{1}{2}\alpha t^2 $$
$$ \Delta\theta = (1.2)(10) + \frac{1}{2}(-0.62)(10)^2 $$
$$ \Delta\theta = 12 - 0.31(100) $$
$$ \Delta\theta = 12 - 31 = -19 \, \si{rad} $$

El signo negativo indica que, al cabo de $10 \, \si{s}$, el desplazamiento neto ha sido en sentido horario debido a la desaceleración e inversión del movimiento.
Para hallar el número de revoluciones (como cantidad escalar), se divide la magnitud del desplazamiento angular entre $2\pi$:
$$ \text{Nro. de revoluciones} = \frac{|\Delta\theta|}{2\pi} = \frac{19}{2\pi} \, \text{rev} $$
$$ \text{Nro. de revoluciones} \approx 3.02 \, \text{rev} $$

\begin{tcolorbox}[colback=green!5!white, colframe=green!50!black, title=Respuesta Final]
\centering \Large $3.02 \, \text{rev}$
\end{tcolorbox}

\newpage

\subsection*{Enunciado}
Una partícula se desplaza con MCUVR en el plano $xy$, al instante $t = 0 \, \si{s}$ pasa por el punto $A(4; 3) \, \si{m}$, con rapidez de $6 \, \si{m/s}$, en sentido antihorario y una aceleración tangencial de $3.1 \, \si{m/s^2}$. Luego de $1.5 \, \si{s}$, pasa por primera vez por el punto $B(-5; 0) \, \si{m}$. Determine su aceleración centrípeta al instante $1.5 \, \si{s}$.

\subsection*{Solución}

\subsubsection*{1. Velocidad angular a los 1.5 s}
Del análisis previo se conocen la velocidad angular inicial ($\omega_0 = 1.2 \, \si{rad/s}$) y la aceleración angular ($\alpha = -0.62 \, \si{rad/s^2}$).
Se calcula la velocidad angular final ($\omega_f$) a los $1.5 \, \si{s}$:
$$ \omega_f = \omega_0 + \alpha t $$
$$ \omega_f = 1.2 + (-0.62)(1.5) $$
$$ \omega_f = 1.2 - 0.93 = 0.27 \, \si{rad/s} $$

En forma vectorial, el movimiento a los $1.5 \, \si{s}$ sigue siendo antihorario:
$$ \vec{\omega}_f = 0.27\vec{k} \, \si{rad/s} $$

\subsubsection*{2. Vector posición a los 1.5 s}
El enunciado indica que a los $1.5 \, \si{s}$ la partícula pasa por el punto $B(-5; 0) \, \si{m}$. 
Por lo tanto, el vector posición de la partícula en ese instante es:
$$ \vec{r}_B = -5\vec{i} \, \si{m} $$

\subsubsection*{3. Cálculo de la aceleración centrípeta}
La aceleración centrípeta ($\vec{a}_c$) se dirige hacia el centro de la trayectoria circular (el origen en este caso). Se calcula mediante la ecuación vectorial:
$$ \vec{a}_c = -\omega_f^2 \vec{r}_B $$
$$ \vec{a}_c = -(0.27)^2 (-5\vec{i}) $$
$$ \vec{a}_c = -0.0729 (-5\vec{i}) $$
$$ \vec{a}_c = 0.3645\vec{i} \, \si{m/s^2} $$

Aproximando a tres cifras decimales, se obtiene el resultado final.

\begin{tcolorbox}[colback=green!5!white, colframe=green!50!black, title=Respuesta Final]
\centering \Large $\vec{a}_c = 0.365\vec{i} \, \si{m/s^2}$
\end{tcolorbox}

\newpage

\subsection*{Enunciado}
Una partícula que describe una circunferencia de radio $10 \, \si{m}$ en el plano $xy$ con MCUV en sentido antihorario. A los $0 \, \si{s}$, su velocidad es $\vec{v}_0 = 3\vec{i} - 4\vec{j} \, \si{m/s}$, que forma un ángulo de $120^\circ$ con su aceleración. Realice el gráfico velocidad angular del radio vector versus tiempo de $0$ a $10 \, \si{s}$.

\subsection*{Solución}

\subsubsection*{1. Análisis de aceleración}
Se calcula la rapidez inicial ($v_0$):
$$ v_0 = \sqrt{3^2 + (-4)^2} = 5 \, \si{m/s} $$

La velocidad angular inicial ($\omega_0$) para el movimiento antihorario es:
$$ \omega_0 = \frac{v_0}{R} = \frac{5}{10} = 0.5 \, \si{rad/s} $$

La aceleración centrípeta ($a_c$) se dirige hacia el centro, formando $90^\circ$ con la velocidad tangencial:
$$ a_c = \frac{v_0^2}{R} = \frac{5^2}{10} = 2.5 \, \si{m/s^2} $$

El ángulo total entre la velocidad y la aceleración es $120^\circ$. El ángulo entre la aceleración total y la aceleración centrípeta es $120^\circ - 90^\circ = 30^\circ$.
Como el ángulo supera los $90^\circ$, la aceleración tangencial ($a_t$) se opone a la velocidad (movimiento retardado):
$$ \tan(30^\circ) = \frac{a_t}{a_c} $$
$$ a_t = a_c \tan(30^\circ) = 2.5 \left(\frac{\sqrt{3}}{3}\right) \approx 1.443 \, \si{m/s^2} $$

\subsubsection*{2. Ecuación de la velocidad angular}
Se halla la aceleración angular ($\alpha$). Al ser un movimiento retardado antihorario, $\alpha$ es negativa:
$$ \alpha = -\frac{a_t}{R} = -\frac{1.443}{10} = -0.1443 \, \si{rad/s^2} $$

La ecuación de la velocidad angular en función del tiempo es:
$$ \omega(t) = \omega_0 + \alpha t $$
$$ \omega(t) = 0.5 - 0.1443t $$

Se determina el tiempo en que la partícula se detiene momentáneamente ($\omega = 0$):
$$ 0 = 0.5 - 0.1443t \implies t = \frac{0.5}{0.1443} \approx 3.464 \, \si{s} $$

La velocidad angular a los $10 \, \si{s}$ es:
$$ \omega(10) = 0.5 - 0.1443(10) = -0.943 \, \si{rad/s} $$

\subsubsection*{3. Gráfico $\omega$ vs t}
\begin{lstlisting}
import matplotlib.pyplot as plt
import numpy as np

t = np.linspace(0, 10, 200)
w = 0.5 - 0.1443 * t

plt.figure(figsize=(7, 4))
plt.plot(t, w, 'b-', linewidth=2)
plt.axhline(0, color='black', linewidth=1)
plt.axvline(0, color='black', linewidth=1)

t_fill1 = np.linspace(0, 3.464, 100)
w_fill1 = 0.5 - 0.1443 * t_fill1
plt.fill_between(t_fill1, w_fill1, 0, color='blue', alpha=0.2)

t_fill2 = np.linspace(3.464, 10, 100)
w_fill2 = 0.5 - 0.1443 * t_fill2
plt.fill_between(t_fill2, w_fill2, 0, color='red', alpha=0.2)

plt.xlabel('t (s)')
plt.ylabel(r'$\omega$ (rad/s)')
plt.grid(True, linestyle='--', alpha=0.6)
plt.xlim(0, 10)

plt.text(1.5, 0.15, r'$\Delta\theta > 0$', fontsize=12, color='blue')
plt.text(6.5, -0.4, r'$\Delta\theta < 0$', fontsize=12, color='red')
plt.text(3.464, 0.05, '3.46s', fontsize=10, ha='center')

plt.tight_layout()
plt.show()
\end{lstlisting}

\newpage

\subsection*{Enunciado}
Una partícula que describe una circunferencia de radio $10 \, \si{m}$ en el plano $xy$ con MCUV en sentido antihorario. A los $0 \, \si{s}$, su velocidad es $\vec{v}_0 = 3\vec{i} - 4\vec{j} \, \si{m/s}$, que forma un ángulo de $120^\circ$ con su aceleración. Determine la distancia recorrida por la partícula, entre $0$ y $10 \, \si{s}$. 

\subsection*{Solución}

\subsubsection*{1. Desplazamiento angular por intervalos}
La distancia recorrida es un escalar dependiente de la longitud de arco. Dado que la partícula invierte su sentido a los $t = 3.464 \, \si{s}$, se debe calcular la distancia angular absoluta sumando los desplazamientos de cada intervalo.

Para el primer intervalo ($0$ a $3.464 \, \si{s}$):
$$ \Delta\theta_1 = \left( \frac{\omega_0 + \omega_f}{2} \right) t_1 $$
$$ \Delta\theta_1 = \left( \frac{0.5 + 0}{2} \right) (3.464) = 0.866 \, \si{rad} $$

Para el segundo intervalo ($3.464 \, \si{s}$ a $10 \, \si{s}$), el tiempo transcurrido es $\Delta t_2 = 10 - 3.464 = 6.536 \, \si{s}$:
$$ \Delta\theta_2 = \left( \frac{\omega_i + \omega_f}{2} \right) \Delta t_2 $$
$$ \Delta\theta_2 = \left( \frac{0 + (-0.943)}{2} \right) (6.536) = -3.082 \, \si{rad} $$

\subsubsection*{2. Distancia total recorrida}
La distancia angular total ($\theta_{total}$) es la suma de los valores absolutos de los desplazamientos:
$$ \theta_{total} = |\Delta\theta_1| + |\Delta\theta_2| $$
$$ \theta_{total} = 0.866 + 3.082 = 3.948 \, \si{rad} $$

La distancia lineal ($d$) se obtiene multiplicando por el radio de la trayectoria ($R = 10 \, \si{m}$):
$$ d = \theta_{total} R $$
$$ d = (3.948)(10) $$
$$ d = 39.48 \, \si{m} $$

\textit{Nota: Redondeando el resultado analítico a una cifra decimal, coincide exactamente con la respuesta del libro de $39.5 \, \si{m}$. El $39.4 \, \si{m}$ del texto original se debe al uso de $\alpha$ aproximado a $-0.144 \, \si{rad/s^2}$ en los pasos intermedios.}

\begin{tcolorbox}[colback=green!5!white, colframe=green!50!black, title=Respuesta Final]
\centering \Large $d = 39.49 \, \si{m}$
\end{tcolorbox}

\end{document}
